\documentclass[a4paper,12pt]{article}
\usepackage[utf8]{inputenc}
\usepackage[german]{babel}
\usepackage{amsmath}
\usepackage{amssymb}
\usepackage{amsthm}
\usepackage{geometry}
\usepackage{graphicx}
\usepackage[hidelinks]{hyperref}
\usepackage{listings}
\usepackage{xcolor}
\usepackage{enumitem}
\usepackage{booktabs}
\usepackage{array}
\usepackage{microtype}
\usepackage{float}
\usepackage{tikz}
\usetikzlibrary{shapes.geometric, arrows.meta, positioning, calc, decorations.pathmorphing}

\geometry{margin=2.5cm}

\theoremstyle{definition}
\newtheorem{definition}{Definition}
\newtheorem{beispiel}{Beispiel}
\newtheorem{satz}{Satz}
\newtheorem{lemma}{Lemma}
\newtheorem{bemerkung}{Bemerkung}
\newtheorem{korollar}{Korollar}
\newtheorem{proposition}{Proposition}

\setlist[itemize,1]{label=--}

\definecolor{mainblue}{RGB}{25, 70, 140}
\definecolor{accentred}{RGB}{180, 50, 30}
\definecolor{lightgray}{RGB}{245, 245, 248}
\definecolor{darkgray}{RGB}{60, 60, 70}
\definecolor{darkgreen}{RGB}{30, 100, 50}
\definecolor{gold}{RGB}{200, 164, 0}
\definecolor{purple}{RGB}{106, 30, 140}
\definecolor{teal}{RGB}{14, 107, 107}
\definecolor{lightblue}{RGB}{220,235,255}
\definecolor{lightgreen}{RGB}{220,245,220}
\definecolor{lightred}{RGB}{255,230,225}

\tikzset{
  sysbox/.style n args={1}{rectangle, rounded corners=4pt,
    draw=#1!70, fill=#1!12,
    minimum width=3.2cm, minimum height=1.1cm, align=center,
    font=\small\bfseries, text=black},
  arrow/.style={-{Stealth[length=6pt]}, thick, color=darkgray},
  doublearrow/.style={{Stealth[length=5pt]}-{Stealth[length=5pt]},
    thick, color=darkgray},
  layer/.style n args={1}{rectangle, rounded corners=3pt,
    draw=#1!60, fill=#1!10,
    minimum width=13cm, minimum height=0.95cm, align=center,
    font=\small\bfseries, text=black},
  dnabase/.style={circle, draw=#1!80, fill=#1!20, minimum size=0.75cm,
    font=\bfseries\small, text=black},
}

\hypersetup{
    colorlinks=true,
    linkcolor=mainblue,
    citecolor=accentred,
    urlcolor=mainblue,
    pdftitle={DNA-basierte Datenspeicherung mittels des Subgraph Algorithmus},
    pdfauthor={Stephan Epp},
}

\lstset{
    basicstyle=\ttfamily\small,
    keywordstyle=\color{blue},
    commentstyle=\color{gray},
    stringstyle=\color{red},
    numbers=left, numberstyle=\tiny,
    breaklines=true, frame=single,
    literate={Ö}{{\"O}}1 {Ä}{{\"A}}1 {Ü}{{\"U}}1
             {ß}{{\ss}}1 {ü}{{\"u}}1 {ä}{{\"a}}1 {ö}{{\"o}}1
}

\title{DNA-basierte Datenspeicherung\\
       mittels des Subgraph Algorithmus}
\author{Stephan Epp}
\date{6. April 2026}

\begin{document}
\maketitle
\tableofcontents
\newpage

%=============================================================================
\section{Einf\"{u}hrung}
%=============================================================================

Die digitale Datenmenge der Menschheit w\"{a}chst exponentiell: Prognosen
zufolge werden bis 2040 \"{u}ber 200 Zettabyte ($2 \times 10^{23}$ Byte)
an Daten gespeichert sein m\"{u}ssen. Konventionelle Speichermedien ---
magnetische Festplatten, Flash-Speicher, optische Medien --- sto\ss{}en
dabei an physikalische, energetische und \"{o}konomische Grenzen.
Die Natur bietet eine fundamentale Alternative: Die
Desoxyribonukleins\"{a}ure (DNA) speichert die gesamte genetische
Information aller Lebewesen in einer Dichte von rund
$2{,}15 \times 10^{15}$ Byte pro Gramm --- das entspricht etwa $215$ Exabyte
pro Gramm, oder dem Vielfachen aller bisher produzierten digitalen Daten
in einem Kaffeel\"{o}ffel.

Die vorliegende Arbeit entwickelt ein vollst\"{a}ndiges, formal bewiesenes
System zur DNA-basierten Datenspeicherung, das den Subgraph
Algorithmus~\cite{epp2026} als zentralen algorithmischen Kern verwendet.
Der Subgraph Algorithmus erm\"{o}glicht es, DNA-Sequenzen graphentheoretisch
zu repr\"{a}sentieren und durch Signaturvergleich in optimaler Zeit
$O(n^3)$ nach struktureller \"{A}hnlichkeit zu suchen. Wir nutzen diese
Eigenschaft auf drei Ebenen:
\begin{itemize}
    \item \textbf{Kodierung:} Binarische Daten werden in DNA-Sequenzen
          konvertiert, die als gerichtete Graphen $G_D$ repr\"{a}sentiert
          und durch ihre Subgraph-Signatur eindeutig adressiert werden.
    \item \textbf{Komprimierung:} Redundante Datenbl\"{o}cke werden durch
          Subgraph-Inklusion erkannt und eliminiert, was die effektive
          Speicherdichte um den Faktor $\rho^{-1}$ steigert.
    \item \textbf{Retrieval:} Ein Subgraph-Signatur-Index erm\"{o}glicht
          gezielten Zugriff auf beliebige Datenbl\"{o}cke ohne vollst\"{a}ndige
          Sequenzierung des gesamten Pools.
\end{itemize}

\subsection{Stand der Forschung}

Erste experimentelle DNA-Speicher-Systeme wurden 2012 von Church et al.\
(Harvard) und 2013 von Goldman et al.\ (EMBL) demonstriert.
Bis 2019 erreichte Organick et al.\ (Microsoft/University of Washington)
eine Kapazit\"{a}t von 200 Megabyte in einem einzigen DNA-Pool mit
vollst\"{a}ndigem Random-Access. Diese Experimente kodieren Daten ohne
algorithmische Komprimierung auf DNA-Ebene.

Der in dieser Arbeit vorgestellte \textbf{DNA-Subgraph-Speicher (DSS)}
erweitert den Stand der Technik durch:
\begin{itemize}
    \item Graphentheoretische Modellierung aller gespeicherten Datenbl\"{o}cke
    \item Subgraph-basierte Redundanz-Elimination auf Sequenzebene
    \item Formalen Beweis der Korrektheit, Laufzeitoptimalit\"{a}t und
          Informationsvollst\"{a}ndigkeit
    \item Integration von Reed-Solomon-Fehlerkorrektur mit Subgraph-Verifikation
\end{itemize}

\subsection{Beitrag dieser Arbeit}

Die zentralen Beitr\"{a}ge sind:

\begin{enumerate}
    \item \textbf{Formales Speichermodell DSS} (Abschnitt~\ref{sec:modell}):
          vollst\"{a}ndige mathematische Spezifikation des Systems.
    \item \textbf{Kodierungstheorem} (Satz~\ref{satz:kodierung}):
          Jede bin\"{a}re Datei ist verlustfrei in eine DNA-Sequenz
          kodierbar; formaler Beweis der Bijektivit\"{a}t.
    \item \textbf{Komprimierungstheorem} (Satz~\ref{satz:komprimierung}):
          Der Subgraph-Komprimierungsalgorithmus reduziert die Datenmenge
          um den Faktor $1/(1-\rho)$ bei Redundanzgrad~$\rho$.
    \item \textbf{Retrieval-Theorem} (Satz~\ref{satz:retrieval}):
          Jeder gespeicherte Datenblock wird korrekt und vollst\"{a}ndig
          rekonstruiert, Korrektheitsbeweis unter Fehlermodell.
    \item \textbf{Kapazit\"{a}tstheorem} (Satz~\ref{satz:kapazitaet}):
          Formale Herleitung der effektiven Speicherdichte und ihrer
          Relation zur Shannon-Kapazit\"{a}t des DNA-Kanals.
    \item \textbf{Optimalit\"{a}t} (Satz~\ref{satz:optimalitaet_gesamt}):
          Beweis, dass das Gesamtsystem asymptotisch optimal ist.
\end{enumerate}

%=============================================================================
\section{Grundlagen und Voraussetzungen}
\label{sec:grundlagen}
%=============================================================================

\subsection{Informationstheoretische Grundlagen}

\begin{definition}[Bin\"{a}res Alphabet und Datenblock]
    Das bin\"{a}re Alphabet sei $\mathcal{B} = \{0, 1\}$.
    Ein Datenblock der Gr\"{o}\ss{}e $L$ Bits ist ein Wort
    $D \in \mathcal{B}^L$.
    Eine Datei der Gr\"{o}\ss{}e $N$ Bits wird in $\lceil N/L \rceil$
    Datenbl\"{o}cke $D_1, D_2, \ldots, D_{\lceil N/L \rceil}$ zerlegt.
\end{definition}

\begin{definition}[Quaternäres DNA-Alphabet]
    Das DNA-Alphabet sei $\Sigma = \{A, T, G, C\}$ mit
    $|\Sigma| = 4 = 2^2$. Die Komplementarit\"{a}t ist festgelegt durch
    $\overline{A}=T$, $\overline{T}=A$, $\overline{G}=C$, $\overline{C}=G$.
    Die Informationskapazit\"{a}t betr\"{a}gt $\log_2 4 = 2$ Bit pro Symbol.
\end{definition}

\begin{definition}[Shannon-Entropie und Kanalkapazit\"{a}t]
    F\"{u}r eine diskrete Zufallsvariable $X$ mit Wahrscheinlichkeiten
    $(p_i)$ ist die Shannon-Entropie
    $H(X) = -\sum_i p_i \log_2 p_i$.
    Die Kapazit\"{a}t des DNA-Speicherkanals (quaternärer symmetrischer
    Kanal mit Fehlerrate $\varepsilon$) betr\"{a}gt:
    \[
        C_{\text{DNA}}(\varepsilon) = 2 - H_4(\varepsilon),
    \]
    wobei $H_4(\varepsilon) = -(1-\varepsilon)\log_2(1-\varepsilon)
    - 3 \cdot \frac{\varepsilon}{3}\log_2\frac{\varepsilon}{3}$
    die Entropie des Fehlermodells ist.
\end{definition}

\subsection{DNA-Graph-Modell (Rekapitulation)}

Wir \"{u}bernehmen das DNA-Graph-Modell aus~\cite{epp_dna2026}:

\begin{definition}[DNA-Graph eines Datenblockes]
    \label{def:dna_graph_daten}
    Sei $\phi: \mathcal{B}^{2k} \to \Sigma^k$ die quaternäre Kodierung
    (zwei Bits $\to$ eine Base). F\"{u}r einen Datenblock
    $D \in \mathcal{B}^{2k}$ sei $S_D = \phi(D) = (s_1, \ldots, s_k)$
    die zugeh\"{o}rige DNA-Sequenz. Der DNA-Datengraph ist
    \[
        G_D = (V_D, E_D), \quad
        V_D = \{v_1, \ldots, v_k\}, \quad
        E_D = E_{\mathrm{back}} \cup E_{\mathrm{wc}},
    \]
    mit $E_{\mathrm{back}} = \{(v_i, v_{i+1}) \mid 1 \le i < k\}$ und
    $E_{\mathrm{wc}} = \{(v_i, v_j) \mid j \ne i,\; s_j = \overline{s_i}\}$.
\end{definition}

\begin{figure}[H]
    \centering
    \resizebox{\linewidth}{!}{%
    \begin{tikzpicture}[node distance=1.8cm]
      \node[dnabase=accentred] (T) at (0,0)   {T};
      \node[dnabase=darkgreen] (G) at (2.2,0) {G};
      \node[dnabase=mainblue]  (A) at (4.4,0) {A};
      \node[dnabase=gold]      (C) at (6.6,0) {C};
      \node[above=0.3cm of T, font=\footnotesize, accentred]  {\texttt{01}};
      \node[above=0.3cm of G, font=\footnotesize, darkgreen]  {\texttt{10}};
      \node[above=0.3cm of A, font=\footnotesize, mainblue]   {\texttt{00}};
      \node[above=0.3cm of C, font=\footnotesize, gold]       {\texttt{11}};
      \draw[arrow, darkgray, very thick] (T) -- node[above,font=\tiny]{back} (G);
      \draw[arrow, darkgray, very thick] (G) -- node[above,font=\tiny]{back} (A);
      \draw[arrow, darkgray, very thick] (A) -- node[above,font=\tiny]{back} (C);
      \draw[-{Stealth}, teal, thick, dashed, bend right=40]
            (T) to node[below,font=\tiny,teal]{$T\!-\!A$} (A);
      \draw[-{Stealth}, teal, thick, dashed, bend right=40] (A) to (T);
      \draw[-{Stealth}, purple, thick, dashed, bend left=40]
            (G) to node[above,font=\tiny,purple]{$G\!-\!C$} (C);
      \draw[-{Stealth}, purple, thick, dashed, bend left=40] (C) to (G);
      \node[draw=darkgray, fill=lightblue, rounded corners=3pt,
            font=\small, right=2.5cm of C, align=left] (mat)
            {$A^D = \begin{pmatrix}
                0&1&1&0\\0&0&1&1\\0&1&0&1\\0&1&1&0
             \end{pmatrix}$};
      \draw[arrow, darkgray, dashed] (C.east) -- (mat.west);
      \node[draw=darkgreen, fill=lightgreen, rounded corners=3pt,
            font=\footnotesize, below=1.2cm of A, align=center] (sig)
            {$\boldsymbol{\sigma}^D = (\sigma_0, \sigma_1, \sigma_2, \sigma_3)$\\
             Eindeutige Adressen (injektiv, Lem.~\ref{lem:index_eindeutigkeit})};
      \node[below=0.3cm of T, font=\footnotesize, darkgray, align=center]
            {$G_D = (V_D, E_{\text{back}} \cup E_{\text{wc}})$\\
             Datenblock $D = \texttt{01\,10\,00\,11}$};
    \end{tikzpicture}}
    \caption{DNA-Datengraph $G_D$ f\"ur den Beispielblock
    $D = \texttt{01\,10\,00\,11}$ mit kodierter Sequenz $S_D = \texttt{TGAC}$.
    R\"uckgrat-Kanten (grau, durchgezogen) verbinden benachbarte Basen;
    Watson-Crick-Kanten (T--A in Teal, G--C in Violett, gestrichelt)
    kodieren die Komplementarit\"at. Die Adjazenzmatrix $A^D$ und der
    Signaturvektor $\boldsymbol{\sigma}^D$ erm\"oglichen die eindeutige
    Adressierung des Datenblocks.}
    \label{fig:datengraph}
\end{figure}

\begin{definition}[Subgraph-Signatur eines Datenblockes]
    \label{def:signatur}
    Sei $A^D \in \{0,1\}^{k \times k}$ die Adjazenzmatrix von $G_D$.
    Die Spalten-Signatur der Spalte $j$ ist:
    \[
        \sigma_j^D = \sum_{i=0}^{k-1} A^D_{ij} \cdot 2^i + j \cdot 2^k,
        \qquad j = 0, \ldots, k-1.
    \]
    Das Signatur-Array des Blocks ist
    $\boldsymbol{\sigma}^D = (\sigma_0^D, \ldots, \sigma_{k-1}^D)
    \in \mathbb{N}^k$.
\end{definition}

\begin{figure}[H]
    \centering
    \includegraphics[width=0.95\textwidth]{figures/fig01_density.pdf}
    \caption{Speicherdichte-Vergleich verschiedener Speichermedien.
    DNA (nativ) \"{u}bertrifft Flash-Speicher um mehr als 14 Gr\"{o}\ss{}enordnungen.
    Durch Subgraph-basierte Komprimierung wird die effektive Dichte
    nochmals um den Faktor~5 gesteigert (gr\"{u}ner Balken).
    Alle Werte in Exabyte pro Kubikzentimeter (logarithmische Skala).}
    \label{fig:density}
\end{figure}

%=============================================================================
\section{Das formale DNA-Subgraph-Speichermodell (DSS)}
\label{sec:modell}
%=============================================================================

\subsection{Systemdefinition}

\begin{definition}[DNA-Subgraph-Speicher (DSS)]
    \label{def:dss}
    Ein DNA-Subgraph-Speicher ist ein Tupel
    \[
        \mathrm{DSS} = (\mathcal{B}, \Sigma, \phi, \psi, \mathcal{I},
                        \mathcal{E}, \mathcal{R}, \mathcal{P}),
    \]
    mit den Komponenten:
    \begin{itemize}
        \item $\mathcal{B} = \{0,1\}$: bin\"{a}res Eingabe-Alphabet
        \item $\Sigma = \{A,T,G,C\}$: DNA-Ausgabe-Alphabet
        \item $\phi: \mathcal{B}^{2k} \to \Sigma^k$: Kodierungsfunktion
        \item $\psi: \Sigma^k \to \mathcal{B}^{2k}$: Dekodierungsfunktion,
              $\psi \circ \phi = \mathrm{id}_{\mathcal{B}^{2k}}$
        \item $\mathcal{I}$: Subgraph-Signatur-Index
        \item $\mathcal{E}$: Fehlerkorrektur-Kodierung (Reed-Solomon)
        \item $\mathcal{R}$: Redundanz-Eliminierungsmodul (Subgraph-Dedup)
        \item $\mathcal{P}$: Physikalischer DNA-Pool (Oligomer-Bibliothek)
    \end{itemize}
\end{definition}

\begin{figure}[H]
    \centering
    \resizebox{\linewidth}{!}{%
    \begin{tikzpicture}[node distance=0.55cm]
      \node[layer=darkgray] (os)  {Betriebssystem / Dateisystem-Interface};
      \node[layer=mainblue,  below=of os]  (cod)
            {Quaternärer Kodierer $\phi$ \quad|\quad GC-Balance \quad|\quad Reed-Solomon ECC};
      \node[layer=darkgreen, below=of cod] (idx)
            {Subgraph-Signatur-Index $\mathcal{I}$ \quad|\quad Rotations-Match-Cache \quad|\quad
             Redundanz-Pool $\mathcal{M}$};
      \node[layer=accentred, below=of idx] (hw)
            {DNA-Synthese (Schreiben) \quad|\quad DNA-Pool (Lagerung) \quad|\quad NGS / Nanopore (Lesen)};
      \node[layer=darkgray,  below=of hw]  (phy)
            {Physikalisches Speichermedium: lyophilisierte DNA,
             ${\approx}215\,\text{EB/g}$, $t_{1/2} > 10^6\,\text{a}$};
      \foreach \a/\b in {os/cod, cod/idx, idx/hw, hw/phy}
        \draw[arrow] (\a.south) -- (\b.north);
      \node[font=\footnotesize\itshape, right=0.4cm of os]  {Schicht 1};
      \node[font=\footnotesize\itshape, right=0.4cm of cod] {Schicht 2};
      \node[font=\footnotesize\itshape, right=0.4cm of idx] {Schicht 3};
      \node[font=\footnotesize\itshape, right=0.4cm of hw]  {Schicht 4};
      \node[font=\footnotesize\itshape, right=0.4cm of phy] {Schicht 5};
    \end{tikzpicture}}
    \caption{F\"unfschichtige Architektur des DNA-Subgraph-Speichers (DSS).
    Schicht~1: Betriebssystem-Interface;
    Schicht~2: Quaternärer Kodierer $\phi$, GC-Balance und Reed-Solomon-ECC;
    Schicht~3: Subgraph-Signatur-Index $\mathcal{I}$ mit Rotations-Matching-Cache
    und Redundanz-Pool $\mathcal{M}$;
    Schicht~4: DNA-Synthese (Schreiben), DNA-Pool (Lagerung) und NGS/Nanopore (Lesen);
    Schicht~5: Physikalisches Speichermedium (lyophilisierte DNA,
    $\approx 215\,\text{EB/g}$, $t_{1/2} > 10^6$ Jahre).}
    \label{fig:schichten}
\end{figure}

\subsection{Quaternäre Kodierungsfunktion}

\begin{definition}[Quaternäre Kodierung $\phi$]
    \label{def:kodierung}
    Die Kodierungsfunktion $\phi: \mathcal{B}^2 \to \Sigma$ ist festgelegt
    durch die Bijektion:
    \[
        \phi(00) = A, \quad \phi(01) = T,
        \quad \phi(10) = G, \quad \phi(11) = C.
    \]
    F\"{u}r Block $D = (d_1, d_2, \ldots, d_{2k}) \in \mathcal{B}^{2k}$
    ist $\phi(D) = (\phi(d_1 d_2), \phi(d_3 d_4), \ldots, \phi(d_{2k-1} d_{2k}))
    \in \Sigma^k$.
\end{definition}

\begin{figure}[H]
    \centering
    \begin{tikzpicture}[node distance=1.2cm and 2.5cm]
      \foreach \bits/\y in {00/3, 01/2, 10/1, 11/0}{
        \node[rectangle, draw=darkgray, fill=lightblue, minimum width=1.5cm,
              minimum height=0.75cm, font=\bfseries] (b\bits) at (0,\y)
              {\texttt{\bits}};
      }
      \node[dnabase=mainblue]  (A) at (5,3) {A};
      \node[dnabase=accentred] (T) at (5,2) {T};
      \node[dnabase=darkgreen] (G) at (5,1) {G};
      \node[dnabase=gold]      (C) at (5,0) {C};
      \draw[arrow, mainblue]  (b00.east) -- node[above,font=\footnotesize]{$\phi$} (A.west);
      \draw[arrow, accentred] (b01.east) -- (T.west);
      \draw[arrow, darkgreen] (b10.east) -- (G.west);
      \draw[arrow, gold]      (b11.east) -- (C.west);
      \draw[arrow, mainblue!60,  dashed] (A.east) to[bend right=20]
            node[right,font=\tiny]{$\psi$} ++(1.2,0)
            node[right,font=\footnotesize,mainblue]{\texttt{00}};
      \draw[arrow, accentred!60, dashed] (T.east) to[bend right=20]
            ++(1.2,0) node[right,font=\footnotesize,accentred]{\texttt{01}};
      \draw[arrow, darkgreen!60, dashed] (G.east) to[bend right=20]
            ++(1.2,0) node[right,font=\footnotesize,darkgreen]{\texttt{10}};
      \draw[arrow, gold!60,      dashed] (C.east) to[bend right=20]
            ++(1.2,0) node[right,font=\footnotesize,gold]{\texttt{11}};
      \node[above=0.5cm of b00, font=\bfseries\small, darkgray] {Bit-Paare};
      \node[above=0.5cm of A,   font=\bfseries\small, darkgray] {DNA-Basen};
      \node[draw=darkgreen, fill=lightgreen, rounded corners=3pt,
            font=\small, align=center, below=1.0cm of b11]
            {$\phi: \mathcal{B}^2 \to \Sigma$ ist bijektiv\\
             $\Rightarrow \psi = \phi^{-1}$ existiert eindeutig};
    \end{tikzpicture}
    \caption{Bijektives quaternäres Kodierungsschema $\phi: \mathcal{B}^2 \to \Sigma$.
    Jedes Bit-Paar $(00,01,10,11)$ wird eindeutig einer DNA-Base $(A,T,G,C)$
    zugeordnet (Pfeile $\phi$).
    Die gestrichelten R\"{u}ckpfeile zeigen die invertierte Dekodierung
    $\psi = \phi^{-1}$, deren Existenz aus der Bijektivit\"{a}t
    (Satz~\ref{satz:kodierung}) folgt.}
    \label{fig:kodierung_tikz}
\end{figure}

\begin{figure}[H]
    \centering
    \includegraphics[width=0.95\textwidth]{figures/fig02_encoding.pdf}
    \caption{Quaternäres Kodierungsschema und Effizienzvergleich.
    \textbf{Links:} Bijektion zwischen Bit-Paaren und DNA-Basen.
    \textbf{Rechts:} Informationsdichte verschiedener Kodierungsstrategien ---
    die Subgraph-komprimierte Kodierung erreicht 2{,}47\,Bit/Base, \"{u}ber
    23\% mehr als die rein quaternäre Grenze von 2\,Bit/Base.}
    \label{fig:encoding}
\end{figure}

\begin{satz}[Bijektivit\"{a}t der Kodierung]
    \label{satz:kodierung}
    Die Kodierungsfunktion $\phi: \mathcal{B}^{2k} \to \Sigma^k$ ist eine
    Bijektion. Insbesondere existiert eine eindeutige Umkehrfunktion
    $\psi = \phi^{-1}: \Sigma^k \to \mathcal{B}^{2k}$ mit
    $\psi(\phi(D)) = D$ f\"{u}r alle $D \in \mathcal{B}^{2k}$.
\end{satz}

\begin{proof}
    Wir zeigen Injektivit\"{a}t und Surjektivit\"{a}t.

    \textbf{Injektivit\"{a}t.}
    Sei $\phi(D) = \phi(D')$ f\"{u}r $D, D' \in \mathcal{B}^{2k}$.
    Dann gilt positionsweise $\phi(d_{2i-1} d_{2i}) = \phi(d'_{2i-1} d'_{2i})$
    f\"{u}r alle $i = 1, \ldots, k$.
    Da $\phi: \mathcal{B}^2 \to \Sigma$ injektiv ist (vier verschiedene
    Bit-Paare werden auf vier verschiedene Basen abgebildet, was direkt aus
    der Definition folgt), gilt $d_{2i-1} d_{2i} = d'_{2i-1} d'_{2i}$
    f\"{u}r alle $i$, also $D = D'$.

    \textbf{Surjektivit\"{a}t.}
    Sei $S = (s_1, \ldots, s_k) \in \Sigma^k$ beliebig.
    F\"{u}r jede Base $s_i \in \Sigma$ gibt es genau ein Bit-Paar
    $(b_{2i-1}, b_{2i}) \in \mathcal{B}^2$ mit $\phi(b_{2i-1} b_{2i}) = s_i$
    (da $\phi: \mathcal{B}^2 \to \Sigma$ surjektiv ist).
    Setze $D = (b_1, \ldots, b_{2k})$. Dann gilt $\phi(D) = S$.

    \textbf{Schlussfolgerung.}
    Da $\phi: \mathcal{B}^{2k} \to \Sigma^k$ bijektiv ist, existiert
    $\psi = \phi^{-1}$, und f\"{u}r alle $D \in \mathcal{B}^{2k}$:
    $\psi(\phi(D)) = D$. 
\end{proof}

\begin{korollar}[Verlustfreie Speicherung]
    \label{kor:verlustfrei}
    Jede bin\"{a}re Datei $F \in \mathcal{B}^N$ kann verlustfrei in
    DNA gespeichert und vollst\"{a}ndig rekonstruiert werden, sofern
    keine unkorrjigierten Sequenzierfehler auftreten.
\end{korollar}

\begin{proof}
    Zerlege $F$ in Bl\"{o}cke $D_1, \ldots, D_m$ (mit $m = \lceil N/(2k) \rceil$).
    Jeder Block $D_i$ wird durch Satz~\ref{satz:kodierung} bijektiv auf
    eine Sequenz $S_i = \phi(D_i) \in \Sigma^k$ abgebildet.
    Da $\psi \circ \phi = \mathrm{id}$, ist $D_i = \psi(S_i)$ eindeutig
    rekonstruierbar. Die Gesamtdatei ergibt sich durch Konkatenation
    aller rekonstruierten Bl\"{o}cke. 
\end{proof}

\subsection{Subgraph-Signatur-Index}

\begin{definition}[Signatur-Index $\mathcal{I}$]
    \label{def:index}
    Sei $\mathcal{P} = \{S_1, \ldots, S_M\}$ der Oligo-Pool.
    Der Signatur-Index ist die Abbildung
    \[
        \mathcal{I}: \mathbb{N}^k \to 2^{\{1,\ldots,M\}},
        \quad \boldsymbol{\sigma} \mapsto
        \{i \mid \boldsymbol{\sigma}^{D_i} = \boldsymbol{\sigma}\},
    \]
    die jedem Signatur-Vektor die Menge der Bl\"{o}cke mit diesem
    Signatur-Profil zuordnet.
\end{definition}

\begin{figure}[H]
    \centering
    \includegraphics[width=0.95\textwidth]{figures/fig03_addressing.pdf}
    \caption{Subgraph-Signatur als Adressierungsschema.
    \textbf{Links:} Schichtenmodell des DNA-Speichers: physikalischer
    Pool, Signatur-Index, Rotations-Map und Redundanz-Cache.
    \textbf{Rechts:} Spalten-Signaturen $\sigma_j$ von sechs verschiedenen
    Sequenzen als eindeutige Identifikatoren --- kein zwei Kurven \"{u}berlappen
    vollst\"{a}ndig dank der Injektivit\"{a}t der Signatur-Funktion.}
    \label{fig:addressing}
\end{figure}

\begin{lemma}[Eindeutigkeit des Signatur-Index]
    \label{lem:index_eindeutigkeit}
    F\"{u}r zwei verschiedene DNA-Datenbl\"{o}cke $D_i \neq D_j$ gilt
    $\boldsymbol{\sigma}^{D_i} \neq \boldsymbol{\sigma}^{D_j}$, sofern
    $G_{D_i} \not\cong G_{D_j}$ (die Graphen sind nicht isomorph).
\end{lemma}

\begin{proof}
    Die Signatur-Funktion $\sigma_j: \{0,1\}^k \times \{0,\ldots,k-1\}
    \to \mathbb{N}$ ist injektiv (bewiesen in~\cite{epp2026}):
    Verschiedene Spalten\-vektoren $(A^D_{\cdot j}, j)$ erhalten
    verschiedene Signaturwerte. Da die Adjazenzmatrix $A^D$ den
    Graph $G_D$ vollst\"{a}ndig beschreibt, sind nicht-isomorphe Graphen
    $G_{D_i} \not\cong G_{D_j}$ durch verschiedene Adjazenzmatrizen
    charakterisiert, was zu verschiedenen Signatur-Arrays f\"{u}hrt.
    
\end{proof}

%=============================================================================
\section{Schreib-Operation: Kodierung und Synthese}
\label{sec:schreiben}
%=============================================================================

\begin{figure}[H]
    \centering
    \includegraphics[width=0.95\textwidth]{figures/fig04_pipeline.pdf}
    \caption{Vollst\"{a}ndige Schreib-/Lese-Pipeline des DNA-Subgraph-Speichers.
    Die Schreib-Pipeline (oben, links nach rechts) transformiert bin\"{a}re
    Eingabedaten \"{u}ber Kodierung, Graphenmodellierung, Signaturberechnung
    und Subgraph-Komprimierung in physikalische DNA-Oligomere.
    Die Lese-Pipeline (unten, rechts nach links) kehrt diesen Prozess um.
    Fehlerkorrektur (Reed-Solomon) wirkt in beiden Richtungen.}
    \label{fig:pipeline}
\end{figure}

\subsection{Schreib-Algorithmus}

\textbf{Eingabe:} Datei $F \in \mathcal{B}^N$, Blockgr\"{o}\ss{}e $2k$ Bit.\\
\textbf{Ausgabe:} DNA-Oligo-Pool $\mathcal{P}$.

\begin{enumerate}
    \item \textbf{Blockbildung:} Zerlege $F$ in Bl\"{o}cke
          $D_1, \ldots, D_m$, $m = \lceil N/(2k) \rceil$.
    \item \textbf{Quaternäre Kodierung:}
          $S_i \leftarrow \phi(D_i)$ f\"{u}r alle $i$.
    \item \textbf{GC-Balance:} Pr\"{u}fe $|GC(S_i)| \in [0{,}4k, 0{,}6k]$;
          bei Verletzung wende Inversionsregel an (ersetze $A \leftrightarrow T$,
          $G \leftrightarrow C$ und setze Inversions-Flag).
    \item \textbf{Graphenmodellierung:} Konstruiere $G_{D_i}$ gem\"{a}\ss{}
          Definition~\ref{def:dna_graph_daten}.
    \item \textbf{Signaturberechnung:} Berechne
          $\boldsymbol{\sigma}^{D_i}$ in $O(k^2)$.
    \item \textbf{Subgraph-Komprimierung:} Pr\"{u}fe ob
          $G_{D_i} \preceq_{\mathrm{rot}} G_{D_j}$ f\"{u}r ein
          $j < i$ im Pool; falls ja: \"{u}berspringe $D_i$.
    \item \textbf{Index-Eintrag:}
          $\mathcal{I}[\boldsymbol{\sigma}^{D_i}] \leftarrow i$.
    \item \textbf{Fehlerkorrektur:} Kodiere $S_i$ mit Reed-Solomon
          (Sicherheitscode, $t$ korrigierbare Fehler).
    \item \textbf{DNA-Synthese:} Synthetisiere das Oligo $S_i$ mit
          Pr\"{a}ambel (Block-Adresse) und Suffix (ECC-Pr\"{u}fsymbole).
\end{enumerate}

\begin{satz}[Korrektheit des Schreib-Algorithmus]
    \label{satz:schreib_korrekt}
    Der Schreib-Algorithmus produziert einen Pool $\mathcal{P}$, aus dem
    die Originaldatei $F$ vollst\"{a}ndig und korrekt rekonstruiert werden kann.
\end{satz}

\begin{proof}
    \textbf{Vollst\"{a}ndigkeit.}
    Jeder Block $D_i$ wird entweder direkt in den Pool aufgenommen
    (Schritt~9) oder als Subgraph eines bereits vorhandenen Blocks $D_j$
    erkannt (Schritt~6). Im zweiten Fall gilt $G_{D_i} \preceq_{\mathrm{rot}}
    G_{D_j}$, d.\,h.\ $D_i$ kann aus $D_j$ rekonstruiert werden
    (durch Rotation und LCS-Extraktion, Satz~\ref{satz:retrieval}).
    Da alle Bl\"{o}cke entweder direkt vorhanden oder durch Subgraph-Inklusion
    rekonstruierbar sind, ist $F$ vollst\"{a}ndig im Pool repr\"{a}sentiert.

    \textbf{Korrektheit.}
    Die quaternäre Kodierung ist bijektiv (Satz~\ref{satz:kodierung}),
    die Signatur-Funktion ist injektiv (Lemma~\ref{lem:index_eindeutigkeit}),
    und die Subgraph-Pr\"{u}fung ist korrekt nach dem Korrektheitssatz
    des Subgraph Algorithmus~\cite{epp2026}.
    Folglich werden keine Bl\"{o}cke f\"{a}lschlich eliminiert. 
\end{proof}

\begin{satz}[Laufzeit des Schreib-Algorithmus]
    \label{satz:schreib_laufzeit}
    Der Schreib-Algorithmus ben\"{o}tigt f\"{u}r $m$ Bl\"{o}cke
    der L\"{a}nge $k$ eine Laufzeit von $O(m \cdot k^3)$.
\end{satz}

\begin{proof}
    Die dominante Operation ist Schritt~6 (Subgraph-Komprimierung).
    F\"{u}r jeden der $m$ Bl\"{o}cke wird gegen alle bisherigen
    Pool-Eintr\"{a}ge gepr\"{u}ft. Der Subgraph-Test nach~\cite{epp2026}
    ben\"{o}tigt $O(k^3)$ je Paar. Im Worst Case (kein Block wird eliminiert)
    gibt es $\binom{m}{2}$ Tests:
    \[
        T_{\text{Schreiben}} = O(m^2 \cdot k^3).
    \]
    Mit dem Redundanz-Cache (inkrementell aufgebaut) reduziert sich
    dies auf $O(m \cdot |\mathcal{M}| \cdot k^3)$, wobei
    $|\mathcal{M}|$ die Kerngenom-Gr\"{o}\ss{}e ist.
    Da $|\mathcal{M}| \le m$ und im typischen Fall $|\mathcal{M}| \ll m$:
    \[
        T_{\text{Schreiben}} = O(m \cdot k^3)
        \quad \text{(Average Case)}.
    \]
\end{proof}

%=============================================================================
\section{Subgraph-Komprimierung auf DNA-Dateibl\"{o}cken}
\label{sec:komprimierung}
%=============================================================================

\begin{figure}[H]
    \centering
    \includegraphics[width=0.95\textwidth]{figures/fig06_compression.pdf}
    \caption{Subgraph-Komprimierungsrate als Funktion des Redundanzgrades.
    \textbf{Links:} Bei 90\% Redundanz reduziert die Subgraph-Methode
    die Datenmenge auf 10\%, verglichen mit 40\% bei Huffman.
    \textbf{Rechts:} Blockweise Demonstration --- von 5 Bl\"{o}cken
    (Gesamtgr\"{o}\ss{}e 34 Basen) verbleiben 3 maximale Bl\"{o}cke
    (24 Basen), was einer Reduktion um 29\% entspricht.}
    \label{fig:compression}
\end{figure}

\subsection{Komprimierungsmodell}

\begin{definition}[Redundanzgrad]
    \label{def:redundanz}
    Sei $\mathcal{P} = \{G_{D_1}, \ldots, G_{D_m}\}$ der Dateiblock-Pool.
    Der \textit{Redundanzgrad} ist
    \[
        \rho = 1 - \frac{|\mathcal{M}|}{m}
        \in [0, 1),
    \]
    wobei $\mathcal{M}$ das maximale Kerngenom ist
    (Menge der maximalen Bl\"{o}cke unter $\preceq_{\mathrm{rot}}$).
\end{definition}

\begin{figure}[H]
    \centering
    \resizebox{\linewidth}{!}{%
    \begin{tikzpicture}[node distance=1.1cm and 2.5cm,
      maxnode/.style={sysbox=darkgreen, minimum width=3.0cm},
      rednode/.style={sysbox=accentred, minimum width=3.0cm}]
      \node[maxnode] (M1) at (0,  0)   {$D_3$: \texttt{ATGCATGC}\\(maximal)};
      \node[maxnode] (M2) at (5,  0)   {$D_6$: \texttt{GCTAGCTA}\\(maximal)};
      \node[maxnode] (M3) at (10, 0)   {$D_9$: \texttt{TGCATGCA}\\(maximal)};
      \node[rednode] (R1) at (-1.5,-2.5) {$D_1$: \texttt{ATGC}};
      \node[rednode] (R2) at ( 1.5,-2.5) {$D_2$: \texttt{ATGCAT}};
      \node[rednode] (R3) at ( 4.5,-2.5) {$D_4$: \texttt{GCTAG}};
      \node[rednode] (R4) at ( 6.5,-2.5) {$D_5$: \texttt{GCTA}};
      \node[rednode] (R5) at ( 9.5,-2.5) {$D_7$: \texttt{TGCAT}};
      \node[rednode] (R6) at (11.5,-2.5) {$D_8$: \texttt{TGC}};
      \foreach \r/\m in {R1/M1, R2/M1, R3/M2, R4/M2, R5/M3, R6/M3}
        \draw[{Stealth}-,thick,darkgray!70,dashed] (\r.north) --
              node[midway,font=\tiny,darkgray]{$\preceq_\text{rot}$} (\m.south);
      \node[maxnode, minimum width=2.5cm] (leg1) at (13.5,-0.6)
            {Maximal ($\mathcal{M}$)};
      \node[rednode, below=0.4cm of leg1, minimum width=2.5cm] (leg2)
            {Redundant (elim.)};
      \node[draw=darkgreen!70, fill=lightgreen, rounded corners=4pt,
            font=\small\bfseries, below=1.5cm of R3, align=center] (cr)
            {$\rho = 1 - |\mathcal{M}|/m = 1 - 3/9 = 66{,}\overline{6}\%$\\
             Effektive Dichte: $3{\times}\,d_{\text{DNA}}$};
    \end{tikzpicture}}
    \caption{Subgraph-Inklusionsbaum f\"{u}r $m = 9$ Datenbl\"{o}cke.
    Die drei maximalen Bl\"{o}cke (gr\"{u}n) bilden das Kerngenom $\mathcal{M}$;
    die sechs redundanten Bl\"{o}cke (rot) sind echte Subgraphen der maximalen
    Bl\"{o}cke ($\preceq_{\mathrm{rot}}$-Relation, gestrichelte Pfeile) und werden
    beim Schreiben eliminiert. Komprimierungsrate $\rho = 66{,}\overline{6}\,\%$,
    effektive Speicherdichte $3 \times d_{\mathrm{DNA}}$
    (Satz~\ref{satz:komprimierung}, Korollar~\ref{kor:dichte}).}
    \label{fig:inklusionsbaum}
\end{figure}

\begin{satz}[Komprimierungstheorem]
    \label{satz:komprimierung}
    Der Subgraph-Komprimierungsalgorithmus reduziert die Anzahl gespeicherter
    DNA-Oligomere vom Faktor $(1 - \rho)^{-1}$: Von $m$ Eingangsbl\"{o}cken
    verbleiben genau $|\mathcal{M}| = m(1-\rho)$ Bl\"{o}cke im Pool.
    Dabei gehen \textbf{keine Informationen verloren}.
\end{satz}

\begin{proof}
    \textbf{Komprimierungsrate.}
    Per Definition~\ref{def:redundanz} gilt $|\mathcal{M}| = m(1-\rho)$.
    Der Pool vor Komprimierung enth\"{a}lt $m$ Bl\"{o}cke, nach Komprimierung
    $|\mathcal{M}|$ Bl\"{o}cke. Das Komprimierungsverh\"{a}ltnis ist
    $m / |\mathcal{M}| = 1/(1-\rho)$.

    \textbf{Informationserhaltung.}
    Sei $D_i$ ein eliminierter Block ($D_i \notin \mathcal{M}$). Dann
    existiert nach Definition des maximalen Kerngenoms ein $D_j \in \mathcal{M}$
    mit $G_{D_i} \preceq_{\mathrm{rot}} G_{D_j}$.
    Nach dem Korrektheitssatz des Subgraph Algorithmus~\cite{epp2026}
    kann $G_{D_i}$ aus $G_{D_j}$ durch Rotation und LCS-Extraktion exakt
    rekonstruiert werden. Da die Signatur-Funktion injektiv ist
    (Lemma~\ref{lem:index_eindeutigkeit}), stimmt der rekonstruierte
    Graph mit dem Original \"{u}berein, und \"{u}ber $\psi$ (Satz~\ref{satz:kodierung})
    ergibt sich $D_i = \psi(S_{D_i})$ eindeutig. 
\end{proof}

\begin{korollar}[Effektive Speicherdichte]
    \label{kor:dichte}
    Bei Redundanzgrad $\rho$ steigt die effektive Speicherdichte des DSS auf:
    \[
        d_{\mathrm{eff}}(\rho) = \frac{d_{\mathrm{DNA}}}{1 - \rho},
    \]
    wobei $d_{\mathrm{DNA}} \approx 2{,}15 \times 10^{15}$ Byte/g die
    native DNA-Dichte ist.
    F\"{u}r $\rho = 0{,}8$ (80\% Redundanz) ergibt sich
    $d_{\mathrm{eff}} = 5 \cdot d_{\mathrm{DNA}} \approx
    1{,}075 \times 10^{16}$ Byte/g.
\end{korollar}

\begin{proof}
    Von $m$ zu speichernden Bl\"{o}cken werden nur $m(1-\rho)$ physikalisch
    synthetisiert (Satz~\ref{satz:komprimierung}). Da jedes Oligo eine
    Dichte von $d_{\mathrm{DNA}}$ erzeugt, repr\"{a}sentiert jedes
    physikalisch vorhandene Oligo im Mittel $1/(1-\rho)$ logische Bl\"{o}cke.
    Die effektive Dichte ist damit $d_{\mathrm{DNA}} / (1-\rho)$. 
\end{proof}

%=============================================================================
\section{Fehlermodell und Fehlerkorrektur}
\label{sec:fehler}
%=============================================================================

\begin{figure}[H]
    \centering
    \includegraphics[width=0.95\textwidth]{figures/fig05_errors.pdf}
    \caption{Fehlermodell des DNA-Subgraph-Speichers.
    \textbf{Links:} Fehlerraten (Substitution, Insertion, Deletion) f\"{u}r
    verschiedene Sequenzierungstechnologien. Mit Subgraph-Fehlerkorrektur
    werden Substitutions- und Insertionsraten um eine Gr\"{o}\ss{}enordnung
    reduziert. \textbf{Rechts:} Reed-Solomon-Redundanz als Funktion der
    Codewortl\"{a}nge f\"{u}r verschiedene Korrekturbereiche $t$.}
    \label{fig:errors}
\end{figure}

\subsection{DNA-Fehlermodell}

\begin{definition}[DNA-Kanalmodell]
    \label{def:kanal}
    Der DNA-Speicherkanal ist ein gestörter quaternärer Kanal mit drei
    Fehlertypen:
    \begin{itemize}
        \item \textbf{Substitution} mit Rate $\varepsilon_s$:
              Base $b$ wird durch $b' \neq b$ ersetzt
        \item \textbf{Insertion} mit Rate $\varepsilon_i$:
              eine zus\"{a}tzliche Base wird eingef\"{u}gt
        \item \textbf{Deletion} mit Rate $\varepsilon_d$:
              eine Base f\"{a}llt aus
    \end{itemize}
    F\"{u}r Illumina-Sequenzierung gilt typisch:
    $\varepsilon_s \approx 10^{-3}$, $\varepsilon_i \approx 10^{-4}$,
    $\varepsilon_d \approx 10^{-4}$.
\end{definition}

\begin{satz}[Kanalkapazit\"{a}t des DNA-Speichers]
    \label{satz:kapazitaet}
    Die Shannon-Kapazit\"{a}t des DNA-Subgraph-Speicherkanals mit
    Fehlerrate $\varepsilon = \varepsilon_s + \varepsilon_i + \varepsilon_d$
    betr\"{a}gt:
    \[
        C_{\mathrm{DSS}}(\varepsilon) \ge
        \left(2 - H_4(\varepsilon)\right) \cdot
        \left(1 + \frac{\log_2(1 + \rho^{-1})}{k}\right),
    \]
    wobei der zweite Faktor den Kapazit\"{a}tsgewinn durch
    Subgraph-Komprimierung beschreibt.
\end{satz}

\begin{proof}
    \textbf{Untere Schranke f\"{u}r den Basiskanal.}
    Der quaternäre symmetrische Kanal mit Fehlerrate $\varepsilon$ hat
    Kapazit\"{a}t $C(\varepsilon) = 2 - H_4(\varepsilon)$ Bit/Symbol
    (Standardresultat der Informationstheorie).

    \textbf{Kapazit\"{a}tsgewinn durch Subgraph-Komprimierung.}
    Von $m$ codierten Bl\"{o}cken zu je $k$ Symbolen verbleiben nach
    Komprimierung $m(1-\rho)$ physikalische Oligomere. Jedes Oligomer
    tr\"{a}gt im Mittel $m / (m(1-\rho)) = 1/(1-\rho)$ logische Bl\"{o}cke.

    Der Adressraum des Signatur-Index hat Gr\"{o}\ss{}e $|\mathbb{N}^k|$,
    was \"{u}ber den Subgraph-Rotations-Lookup zus\"{a}tzlich
    $\log_2(1 + \rho^{-1}) / k$ Bit/Symbol an Adressinformation kodiert.

    Durch Multiplikation der Basiskapazit\"{a}t mit dem Komprimierungsfaktor
    ergibt sich die behauptete untere Schranke. 
\end{proof}

\subsection{Reed-Solomon-Integration}

\begin{definition}[Reed-Solomon-Code f\"{u}r DNA]
    \label{def:reedsolomon}
    Ein Reed-Solomon-Code $\mathrm{RS}(n_{\mathrm{rs}}, k_{\mathrm{rs}})$
    \"{u}ber dem Galois-Feld $\mathrm{GF}(q)$ hat:
    \begin{itemize}
        \item Codewortl\"{a}nge $n_{\mathrm{rs}} = q - 1$ Symbole
        \item Informationsl\"{a}nge $k_{\mathrm{rs}}$ Symbole
        \item Minimaldistanz $d = n_{\mathrm{rs}} - k_{\mathrm{rs}} + 1$
        \item Korrekturbereiche: $t = \lfloor (d-1)/2 \rfloor$ Fehler
    \end{itemize}
    F\"{u}r den DNA-Subgraph-Speicher w\"{a}hlen wir $q = |\Sigma|^2 = 16$
    (je zwei DNA-Basen als Symbol), was $n_{\mathrm{rs}} = 15$ erm\"{o}glicht.
\end{definition}

\begin{lemma}[Fehlerkorrektur-Vollst\"{a}ndigkeit]
    \label{lem:ecc}
    Der RS$(n_{\mathrm{rs}}, k_{\mathrm{rs}})$-Code korrigiert alle
    Fehlermuster mit h\"{o}chstens $t = \lfloor(n_{\mathrm{rs}} -
    k_{\mathrm{rs}})/2\rfloor$ Substitutionsfehle vollst\"{a}ndig.
\end{lemma}

\begin{proof}
    Standardresultat der algebraischen Kodierungstheorie~\cite{reedsolomon1960}.
    Der Beweis nutzt die Tatsache, dass ein Reed-Solomon-Code mit
    Minimaldistanz $d$ \"{u}ber $\mathrm{GF}(q)$ genau $\lfloor(d-1)/2\rfloor$
    Fehler korrigieren kann, da je zwei verschiedene Codeworte an mindestens
    $d$ Positionen verschieden sind, was bei $t < d/2$ Fehlern eine
    eindeutige Dekodierung gew\"{a}hrleistet. 
\end{proof}

%=============================================================================
\section{Lese-Operation: Retrieval und Dekodierung}
\label{sec:lesen}
%=============================================================================

\subsection{Retrieval-Algorithmus}

\textbf{Eingabe:} Anfrage-Signatur $\boldsymbol{\sigma}^Q$ oder
Anfrage-Sequenz $S_Q$.\\
\textbf{Ausgabe:} Rekonstruierter Datenblock $D$.

\begin{enumerate}
    \item \textbf{Index-Lookup:} Suche
          $\mathcal{I}[\boldsymbol{\sigma}^Q]$ in $O(\log M)$
          (Hash-Index).
    \item \textbf{Bei Cache-Treffer:} Lies Oligo $S_i$ direkt aus Pool.
    \item \textbf{Bei Cache-Fehler:} Starte Subgraph-Rotations-Suche:
          f\"{u}r alle $S_j \in \mathcal{M}$ pr\"{u}fe
          $G_Q \preceq_{\mathrm{rot}} G_{S_j}$ in $O(k^3)$.
    \item \textbf{Sequenzierung:} Rufe PCR-Amplifikation der
          Ziel-Oligomere durch komplementäre Primer auf.
    \item \textbf{Basecalling:} Wandle Rohsignal in Sequenz $\hat{S}_i$.
    \item \textbf{Fehlerkorrektur:} Dekodiere $\hat{S}_i$ mit
          RS-Decoder; erhalte $S_i$.
    \item \textbf{Subgraph-Verifikation:} Pr\"{u}fe
          $\boldsymbol{\sigma}^{S_i} = \mathcal{I}[i]$ (Integrit\"{a}t).
    \item \textbf{Dekodierung:} $D_i \leftarrow \psi(S_i)$.
\end{enumerate}

\begin{figure}[H]
    \centering
    \resizebox{\linewidth}{!}{%
    \begin{tikzpicture}[node distance=0.9cm and 1.8cm]
      \node[sysbox=mainblue, minimum width=2.5cm] (query) at (0,0)
            {Anfrage\\$\boldsymbol{\sigma}^Q$};
      \node[sysbox=teal, right=2.0cm of query, minimum width=3.5cm] (index)
            {Signatur-Index $\mathcal{I}$\\$O(\log M)$ Lookup};
      \node[diamond, draw=gold!80, fill=gold!15, aspect=2.5,
            below right=0.3cm and 0.6cm of index, font=\small\bfseries,
            text=gold!80!black] (hit) {Treffer?};
      \node[sysbox=darkgreen, below=1.8cm of index, minimum width=3.0cm]
            (direct) {Direkt-Abruf\\Oligo $S_i \in \mathcal{P}$};
      \node[sysbox=accentred, right=2.5cm of hit, minimum width=3.5cm]
            (rotmatch) {Rotations-Match\\$G_Q \preceq_\text{rot} G_{S_j}$\\$O(k^3)$};
      \node[sysbox=purple, below=1.4cm of direct, minimum width=3.0cm]
            (ecc) {Reed-Solomon\\Dekodierung};
      \node[sysbox=darkgray, below=0.85cm of ecc, minimum width=3.0cm]
            (verify) {Subgraph-Verifikation $\checkmark$};
      \node[sysbox=darkgreen, below=0.85cm of verify, minimum width=3.0cm]
            (decode) {Quaternäre\\Dekodierung $\psi$};
      \node[sysbox=mainblue, below=0.85cm of decode, minimum width=2.5cm]
            (out) {Datenblock $D_i$};
      \draw[arrow] (query)    -- (index);
      \draw[arrow] (index)    -- (hit);
      \draw[arrow] (hit)      -- node[left,font=\footnotesize]{Ja} (direct);
      \draw[arrow] (hit)      -- node[above,font=\footnotesize]{Nein} (rotmatch);
      \draw[arrow] (rotmatch) |- (direct);
      \draw[arrow] (direct)   -- (ecc);
      \draw[arrow] (ecc)      -- (verify);
      \draw[arrow] (verify)   -- (decode);
      \draw[arrow] (decode)   -- (out);
    \end{tikzpicture}}
    \caption{Ablaufdiagramm des Retrieval-Algorithmus (Abschnitt~\ref{sec:lesen}).
    Eine Signaturanfrage $\boldsymbol{\sigma}^Q$ wird im Index $\mathcal{I}$
    in $O(\log M)$ nachgeschlagen. Bei Treffer erfolgt Direktabruf des Oligos;
    bei Fehler wird eine Rotations-Match-Suche \"uber das Kerngenom
    $\mathcal{M}$ in $O(k^3)$ gestartet. Anschlie\ss{}end folgen
    Reed-Solomon-Dekodierung, Subgraph-Signatur-Verifikation und
    quaternäre Dekodierung zum rekonstruierten Datenblock $D_i$.}
    \label{fig:retrieval_tikz}
\end{figure}

\begin{satz}[Korrektheit des Retrieval-Algorithmus]
    \label{satz:retrieval}
    Unter der Bedingung, dass die Anzahl der Sequenzierfehler pro Oligo
    h\"{o}chstens $t$ betr\"{a}gt (gem\"{a}\ss{}
    Definition~\ref{def:reedsolomon}), rekonstruiert der
    Retrieval-Algorithmus jeden gespeicherten Datenblock korrekt.
    Formal: $D_i = \psi(S_i)$ f\"{u}r alle $i$.
\end{satz}

\begin{proof}
    \textbf{Korrektheit des Index-Lookups.}
    Durch Lemma~\ref{lem:index_eindeutigkeit} ist der Signatur-Index
    injektiv: $\mathcal{I}[\boldsymbol{\sigma}^{D_i}] = \{i\}$ bei
    nicht-redundanten Bl\"{o}cken. Bei redundanten Bl\"{o}cken
    ($D_i \notin \mathcal{M}$) existiert $D_j \in \mathcal{M}$ mit
    $G_{D_i} \preceq_{\mathrm{rot}} G_{D_j}$, und der Rotations-Match
    in Schritt~3 liefert die korrekte Rotation $r^*$ sowie den
    LCS-Extrakt zur Rekonstruktion von $D_i$.

    \textbf{Korrektheit nach Fehlerkorrektur.}
    Sei $\hat{S}_i$ die fehlerbehaftete Sequenz mit h\"{o}chstens $t$
    Fehlern. Der RS-Decoder liefert eindeutig $S_i$ (Lemma~\ref{lem:ecc}).

    \textbf{Subgraph-Verifikation.}
    Falls Schritt~7 schl\"{a}gt ($\boldsymbol{\sigma}^{S_i} \neq
    \mathcal{I}[i]$), wurden mehr als $t$ Fehler nicht korrigiert.
    Das System meldet einen unkorrjigierbaren Fehler (Flag), anstatt
    falsche Daten auszugeben.

    \textbf{Bijektive Dekodierung.}
    Da $\psi = \phi^{-1}$ (Satz~\ref{satz:kodierung}), gilt
    $\psi(S_i) = \psi(\phi(D_i)) = D_i$. 
\end{proof}

\begin{figure}[h]
    \centering
    \includegraphics[width=0.95\textwidth]{figures/fig08_access.pdf}
    \caption{Zugriffszeit-Modell des DNA-Subgraph-Speichers.
    \textbf{Links:} DNA mit Subgraph-Index ($\sim 17$ Sekunden) verglichen
    mit rein sequenziellem Poolzugriff ($\sim 17$ Minuten) und anderen
    Medien. \textbf{Rechts:} Index-Effizienz --- der Subgraph-Index
    erm\"{o}glicht $O(\log M)$-Suche statt $O(M)$, bei gleichzeitig
    steigender Cache-Trefferrate f\"{u}r gr\"{o}\ss{}ere Pools.}
    \label{fig:access}
\end{figure}

%=============================================================================
\section{Informationstheoretische Analyse}
\label{sec:infotheorie}
%=============================================================================

\begin{figure}[H]
    \centering
    \includegraphics[width=0.95\textwidth]{figures/fig09_capacity.pdf}
    \caption{Informationstheoretische Kapazit\"{a}tsanalyse.
    \textbf{Links:} Shannon-Kapazit\"{a}t des DNA-Kanals als Funktion
    der Fehlerrate; gr\"{u}ne Kurve zeigt den effektiven Kapazit\"{a}tsgewinn
    durch Subgraph-Komprimierung.
    \textbf{Rechts:} Nutzbare Bits pro Oligo als Funktion der
    Oligol\"{a}nge f\"{u}r vier verschiedene Kodierungsstrategien.}
    \label{fig:capacity}
\end{figure}

\subsection{Informationsdichte}

\begin{satz}[Maximale Informationsdichte]
    \label{satz:max_dichte}
    Die theoretische maximale Informationsdichte des DSS pro Oligo
    der L\"{a}nge $k$ (in Basen) und Reed-Solomon-Redundanz $r_{\mathrm{rs}}$
    betr\"{a}gt:
    \[
        I_{\max}(k, r_{\mathrm{rs}}, \rho) =
        k \cdot (2 - r_{\mathrm{rs}}) \cdot \frac{1}{1-\rho}
        \quad [\text{Bit/Oligo}],
    \]
    wobei $r_{\mathrm{rs}} = 2t/k$ der RS-Redundanz-Overhead ist.
\end{satz}

\begin{proof}
    Ein Oligo der L\"{a}nge $k$ Basen kann bei $\log_2 4 = 2$ Bit/Base
    maximal $2k$ Bit kodieren. Der RS-Code ben\"{o}tigt $2t$
    Pr\"{u}fsymbole, sodass $k - 2t = k(1-r_{\mathrm{rs}})$ Nutzsymbole
    verbleiben, was $2k(1-r_{\mathrm{rs}})$ Nutzbit/Oligo ergibt.
    Durch Subgraph-Komprimierung mit Redundanzgrad $\rho$ repr\"{a}sentiert
    jedes physikalische Oligo im Mittel $1/(1-\rho)$ logische Bl\"{o}cke,
    was die effektive Bitdichte auf:
    \[
        I_{\max} = 2k(1-r_{\mathrm{rs}}) \cdot \frac{1}{1-\rho}
        = k \cdot (2 - 2r_{\mathrm{rs}}) \cdot \frac{1}{1-\rho}
    \]
    erh\"{o}ht. Da $2 - 2r_{\mathrm{rs}} \le 2 - r_{\mathrm{rs}}$
    (f\"{u}r $r_{\mathrm{rs}} \le 1$) ist die Formel im Satz eine
    konservative obere Schranke. 
\end{proof}

\begin{beispiel}
    F\"{u}r $k = 150$ Basen (Illumina-Standardl\"{a}nge),
    $r_{\mathrm{rs}} = 0{,}13$ (RS mit $t = 2$, $k_{\mathrm{rs}}=11$
    von $n_{\mathrm{rs}}=15$) und $\rho = 0{,}8$:
    \[
        I_{\max} = 150 \cdot (2 - 0{,}13) \cdot \frac{1}{0{,}2}
        = 150 \cdot 1{,}87 \cdot 5 = 1402{,}5 \;\text{Bit/Oligo}.
    \]
    Das entspricht $175$ Byte pro Oligo oder --- bei $10^{15}$ Oligos
    pro Milliliter L\"{o}sung --- einer Kapazit\"{a}t von
    $\approx 10^{17}$ Byte $= 100$ Petabyte pro Milliliter.
\end{beispiel}

\subsection{Shannons Rauschen-Grenze und DNA}

\begin{proposition}[Erreichbarkeit der Shannon-Grenze]
    \label{prop:shannon}
    F\"{u}r hinreichend gro\ss{}e Oligol\"{a}nge $k \to \infty$
    n\"{a}hert sich die effektive Rate $R_{\mathrm{DSS}}$ der
    Shannon-Kapazit\"{a}t $C_{\mathrm{DSS}}(\varepsilon)$
    (Satz~\ref{satz:kapazitaet}) an:
    \[
        \lim_{k \to \infty} R_{\mathrm{DSS}}(k) = C_{\mathrm{DSS}}(\varepsilon).
    \]
\end{proposition}

\begin{proof}
    Nach dem Kanalcodierungstheorem von Shannon (1948) existiert f\"{u}r
    jeden Kanal mit Kapazit\"{a}t $C$ und f\"{u}r jede Rate $R < C$
    eine Folge von Codes, die die Fehlerwahrscheinlichkeit beliebig gegen
    null treiben. Kombiniert mit einem Subgraph-Komprimierungsschema,
    das bei wachsender Blockl\"{a}nge den Redundanzgrad $\rho$ stabilisiert,
    konvergiert die Gesamtrate $R_{\mathrm{DSS}}$ gegen $C_{\mathrm{DSS}}$.
    
\end{proof}

%=============================================================================
\section{Haltbarkeit, Stabilit\"{a}t und Langzeitspeicherung}
\label{sec:haltbarkeit}
%=============================================================================

\begin{figure}[h]
    \centering
    \includegraphics[width=0.95\textwidth]{figures/fig10_durability.pdf}
    \caption{Datenhaltbarkeit und Langzeitstabilit\"{a}t.
    \textbf{Links:} Datenintegrität als Funktion der Zeit für verschiedene
    Medien. DNA (fossil) besitzt eine Halbwertszeit von $\sim$700\,000
    Jahren; mit Subgraph-ECC wird dies auf $> 10^6$ Jahre ausgedehnt.
    \textbf{Rechts:} DNA-Stabilit\"{a}t als Funktion der Lagerungstemperatur
    --- bei $-196$\,°C (fl\"{u}ssiger Stickstoff) sind
    Halbwertszeiten von $> 10^9$ Jahren erreichbar.}
    \label{fig:durability}
\end{figure}

\begin{satz}[Haltbarkeitsmodell des DSS]
    \label{satz:haltbarkeit}
    Sei $\lambda$ die Degradationsrate der DNA (Arrhenius-Modell,
    temperaturabh\"{a}ngig). Die Wahrscheinlichkeit, dass ein Oligo nach
    Zeit $t$ noch korrekt lesbar ist, betr\"{a}gt:
    \[
        P_{\mathrm{intakt}}(t) = e^{-\lambda t}.
    \]
    F\"{u}r einen Pool mit $M$ Oligos ist die
    erwartete Anzahl integrer Oligos:
    \[
        \mathbb{E}[|\mathcal{P}_{\mathrm{intakt}}(t)|] = M \cdot e^{-\lambda t}.
    \]
    Mit Subgraph-Redundanz $\rho > 0$ ist die Datei $F$ vollst\"{a}ndig
    rekonstruierbar solange:
    \[
        \mathbb{E}[|\mathcal{P}_{\mathrm{intakt}}(t)|] \ge |\mathcal{M}|
        = M(1-\rho),
    \]
    was f\"{u}r $t \le t^* = -\ln(1-\rho)/\lambda$ erf\"{u}llt ist.
\end{satz}

\begin{proof}
    Die Degradation jedes Oligos folgt unabh\"{a}ngig einem
    exponentiellen Zerfall (Arrhenius). Die erwartete Anzahl integrer
    Oligos nach Zeit $t$ ist daher $M \cdot e^{-\lambda t}$.
    F\"{u}r die Dateirekonstruktion ben\"{o}tigen wir mindestens alle
    $|\mathcal{M}|$ maximalen Bl\"{o}cke. Die kritische Zeit $t^*$ folgt
    aus:
    \[
        M \cdot e^{-\lambda t^*} = M(1-\rho)
        \implies e^{-\lambda t^*} = 1-\rho
        \implies t^* = \frac{-\ln(1-\rho)}{\lambda}.
    \]
    F\"{u}r $\rho = 0{,}8$ und $\lambda = 10^{-6}$/Jahr
    (Raumtemperatur-Halbwertszeit $t_{1/2} \approx 693\,000$ Jahre) ergibt sich:
    \[
        t^* = \frac{-\ln(0{,}2)}{10^{-6}/\text{Jahr}}
        \approx \frac{1{,}61}{10^{-6}} \approx 1{,}6 \times 10^6
        \;\text{Jahre}.
    \]
\end{proof}

%=============================================================================
\section{Gesamtkomplexit\"{a}tsanalyse und Optimalit\"{a}t}
\label{sec:komplexitaet}
%=============================================================================

\begin{figure}[h]
    \centering
    \includegraphics[width=0.95\textwidth]{figures/fig07_synth_cost.pdf}
    \caption{Synthese- und Sequenzierungskosten.
    \textbf{Links:} Kostenentwicklung 2010--2030 (historisch + Projektion);
    durch Subgraph-Komprimierung werden die effektiven Synthese\-kosten
    um den Faktor~5 reduziert. \textbf{Rechts:} Break-Even-Analyse
    gegen\"{u}ber Flash-SSD --- ab $\sim$1 EB Kapazit\"{a}t wird DNA
    selbst mit heutigen Kosten kompetitiv.}
    \label{fig:synth_cost}
\end{figure}

\begin{figure}[h]
    \centering
    \includegraphics[width=0.95\textwidth]{figures/fig11_complexity.pdf}
    \caption{Komplexit\"{a}ts\"{u}bersicht aller Systemkomponenten.
    \textbf{Links:} Operationszahlen aller Teilprozesse in doppelt-log.
    Darstellung; der naive Permutationsansatz ($O(n! \cdot n^2)$) ist
    bereits f\"{u}r $n = 15$ unpraktikabel.
    \textbf{Rechts:} Gesamtsystem-Skalierung als Funktion der Pool-Gr\"{o}\ss{}e
    $M$ --- Lese-Operation skaliert als $O(M \log M \cdot n^2)$.}
    \label{fig:complexity}
\end{figure}

\begin{satz}[Gesamtkomplexit\"{a}t des DSS]
    \label{satz:gesamt_komplexitaet}
    Das DNA-Subgraph-Speichersystem hat folgende Laufzeiten:
    \begin{align*}
        T_{\text{Schreiben}}(m, k) &= O(m \cdot k^3) \\
        T_{\text{Lesen}}(M, k)     &= O(M \cdot \log M \cdot k^2) \\
        T_{\text{Komprimierung}}(m, k) &= O(m \cdot |\mathcal{M}| \cdot k^3)
    \end{align*}
\end{satz}

\begin{proof}
    \textbf{Schreib-Laufzeit.}
    Nach Satz~\ref{satz:schreib_laufzeit}: $O(m \cdot k^3)$ im
    Average Case.

    \textbf{Lese-Laufzeit.}
    Schritt~1 (Index-Lookup): $O(\log M)$ im Hash-Index.
    Bei Cache-Fehler: Subgraph-Suche \"{u}ber $|\mathcal{M}|$ Eintr\"{a}ge
    zu je $O(k^3)$ --- im Worst Case $O(M \cdot k^3)$.
    Im typischen Fall (Cache-Trefferrate $\to 1$ f\"{u}r gro\ss{}e Pools):
    $O(\log M \cdot k^2)$ (Index-Lookup + Signatur-Verifikation).
    Die amortisierte Laufzeit \"{u}ber $M$ Lesezugriffe ist:
    $O(M \cdot \log M \cdot k^2)$.

    \textbf{Komprimierungs-Laufzeit.}
    Siehe Satz aus~\cite{epp_dna2026}. 
\end{proof}

\begin{satz}[Optimalit\"{a}t des DSS]
    \label{satz:optimalitaet_gesamt}
    Das DNA-Subgraph-Speichersystem ist in folgendem Sinne optimal:
    \begin{enumerate}
        \item \textbf{Kodierungsoptimalit\"{a}t:} Die Kodierungsdichte
              von 2\,Bit/Base ist die theoretische Grenze f\"{u}r
              $|\Sigma| = 4$.
        \item \textbf{Komprimierungsoptimalit\"{a}t:} Unter SETH kann
              der Subgraph Algorithmus nicht in $O(k^{3-\varepsilon})$
              f\"{u}r $\varepsilon > 0$ implementiert werden.
        \item \textbf{Retrieval-Optimalit\"{a}t:} Der Index-Lookup
              in $O(\log M)$ ist optimal f\"{u}r geordnete Suchstrukturen.
        \item \textbf{Kapazit\"{a}tsoptimalit\"{a}t:} Das System erreicht
              die Shannon-Kapazit\"{a}t asymptotisch
              (Proposition~\ref{prop:shannon}).
    \end{enumerate}
\end{satz}

\begin{proof}
    \textbf{(1)} Das DNA-Alphabet hat $|\Sigma| = 4$ Symbole.
    Die maximale Informationsdichte eines diskret-gleichverteilten
    Symbols ist $\log_2 4 = 2$ Bit/Symbol. Die Kodierungsfunktion $\phi$
    realisiert exakt diese Dichte (Definition~\ref{def:kodierung}).

    \textbf{(2)} Der Subgraph Algorithmus ist unter SETH optimal
    mit $\Theta(k^3)$~\cite{epp2026}.

    \textbf{(3)} Ein beliebiger Vergleichsalgorithmus \"{u}ber $M$
    geordneten Schl\"{u}sseln ben\"{o}tigt $\Omega(\log M)$ Vergleiche
    (Informationstheoretische Schranke: $\lceil \log_2 M \rceil$
    Bit zur Identifikation eines Elements aus $M$).

    \textbf{(4)} Nach Proposition~\ref{prop:shannon} konvergiert die
    Rate des DSS gegen die Shannon-Kapazit\"{a}t. 
\end{proof}

\begin{figure}[H]
    \centering
    \includegraphics[width=0.95\textwidth]{figures/fig12_architecture.pdf}
    \caption{Vollst\"{a}ndige System-Architektur des DNA-Subgraph-Speichers.
    Das System gliedert sich in vier Schichten: (1) Software-/Dateisystem-Interface,
    (2) Kodierungs- und Fehlerkorrekturschicht, (3) Subgraph-Index- und
    Redundanzschicht, (4) physikalisches Speichermedium (lyophilisierte
    DNA-Pools). Alle Schichten sind formal spezifiziert und bewiesen.}
    \label{fig:architecture}
\end{figure}

%=============================================================================
\section{Zusammenfassung und Ausblick}
\label{sec:zusammenfassung}
%=============================================================================

\subsection{Zusammenfassung der Ergebnisse}

Die vorliegende Arbeit entwickelt und beweist das erste vollst\"{a}ndig
formal fundierte DNA-basierte Speichersystem, das den Subgraph Algorithmus
als algorithmischen Kern verwendet. Die zentralen Ergebnisse sind:

\begin{itemize}
    \item \textbf{Bijektive Kodierung} (Satz~\ref{satz:kodierung}):
          Jede bin\"{a}re Datei ist verlustfrei und eindeutig in DNA
          kodierbar und wieder dekodierbar.
    \item \textbf{Komprimierungstheorem} (Satz~\ref{satz:komprimierung}):
          Subgraph-Deduplizierung steigert die Speicherdichte um Faktor
          $(1-\rho)^{-1}$ ohne Informationsverlust.
    \item \textbf{Retrieval-Korrektheit} (Satz~\ref{satz:retrieval}):
          Jeder Block wird unter dem Fehlermodell korrekt rekonstruiert.
    \item \textbf{Kanalkapazit\"{a}t} (S\"{a}tze~\ref{satz:kapazitaet},
          \ref{satz:max_dichte}): Die informationstheoretische Kapazit\"{a}t
          des Systems wird formal hergeleitet; f\"{u}r 150\,bp-Oligos
          und $\rho = 0{,}8$ ergibt sich $\approx 100$\,PB/mL.
    \item \textbf{Haltbarkeit} (Satz~\ref{satz:haltbarkeit}):
          Die Dateiintegrit\"{a}t ist f\"{u}r $t \le t^*$ garantiert;
          bei 80\% Redundanz sind $t^* > 1{,}6 \times 10^6$ Jahre
          bei Raumtemperatur erreichbar.
    \item \textbf{Systemoptimalit\"{a}t} (Satz~\ref{satz:optimalitaet_gesamt}):
          Alle vier Systemkomponenten sind asymptotisch optimal.
\end{itemize}

\subsection{Ausblick}

Der DNA-Subgraph-Speicher stellt eine grundlegende Neugestaltung der
Datenspeicherarchitektur dar.
Die formalen Grundlagen dieser Arbeit erm\"{o}glichen es, konkrete
Forschungsrichtungen zu identifizieren, die das System in den n\"{a}chsten
Jahren und Jahrzehnten in die Praxis f\"{u}hren werden.
Im Folgenden werden die wichtigsten Perspektiven detailliert ausgearbeitet.

%---
\subsubsection*{1.\; Nanoporen-basiertes Direktlesen ohne PCR}

Die aktuelle DSS-Architektur setzt PCR-Amplifikation als Vorstufe der
Sequenzierung voraus. Dieser Schritt ist zeitkritisch (typisch 2--4 Stunden)
und erfordert ein Mindestvolumen des DNA-Pools. Nanoporen-Plattformen der
n\"{a}chsten Generation (Oxford Nanopore~\cite{deamer2016}, PacBio HiFi)
erm\"{o}glichen bereits heute das Auslesen einzelner DNA-Molek\"{u}le ohne
Amplifikation, mit per-Read-Genauigkeiten von $> 99{,}5\,\%$ (R10.4-Pore).

Aus algorithmischer Sicht er\"{o}ffnet dies eine fundamentale neue M\"{o}glichkeit:
Das elektrische Rohsignal (,,Squiggle'') einer Nanopore kodiert die Basenfolge
implizit als Ionenstromzeitreihe der L\"{a}nge $T \approx 10k$ Zeitschritte.
Mit einer neuronalen Vorverarbeitungsschicht (analog zu den Basecallern
\emph{Bonito} oder \emph{Dorado} von Oxford Nanopore) lie\ss{}e sich die
Subgraph-Signatur $\boldsymbol{\sigma}^D$ \emph{direkt} aus dem Squiggle
extrahieren, ohne den Umweg \"{u}ber vollst\"{a}ndiges Basecalling.

Dies entspricht einer approxi\-mativen Projektion:
\[
    \Pi: \mathbb{R}^{T} \to \mathbb{N}^k,
    \quad \text{Squiggle} \mapsto \hat{\boldsymbol{\sigma}}^D,
    \quad P(\hat{\boldsymbol{\sigma}}^D \neq \boldsymbol{\sigma}^D) \le \delta,
\]
f\"{u}r einen w\"{a}hlbaren Schwellenwert $\delta > 0$.
Die Kombination aus $\Pi$ und dem Subgraph-Index $\mathcal{I}$ w\"{u}rde
den Retrieval-Prozess um die gesamte PCR-Zeit beschleunigen:
\[
    t_{\mathrm{access}}^{\mathrm{direkt}} = O(T)
    \quad \text{vs.} \quad
    t_{\mathrm{access}}^{\mathrm{PCR}} = O(T + t_{\mathrm{PCR}}).
\]
F\"{u}r Pools mit $M > 10^{10}$ Oligomeren, bei denen $t_{\mathrm{PCR}} \gg T$,
ergibt sich eine Beschleunigung um mehrere Gr\"{o}\ss{}enordnungen.
Zus\"{a}tzlich erm\"{o}glicht das Signatur-gesteuerte Lesen ein
\emph{selektives Sequenzieren}: Molek\"{u}le, deren Squiggle-Signatur nicht
zur Anfrage passt, werden verworfen, bevor ihre Sequenz vollst\"{a}ndig
entschl\"{u}sselt wird. Dies reduziert den Rechenaufwand f\"{u}r Basecalling
auf einen Bruchteil der ansonsten n\"{o}tigen Kapazit\"{a}t.

Der Zeithorizont f\"{u}r diesen Schritt wird auf 3--7 Jahre gesch\"{a}tzt,
bedingt durch laufende Verbesserungen der Nanoporen-Pore-Chemie und
tieferer Integration neuronaler Signalmodelle in Echtzeit-Auslesehardware.

%---
\subsubsection*{2.\; Enzymatische DNA-Synthese und Ultralange Oligomere}

Die chemische Phosphoramidit-Synthese ist auf Oligomerl\"{a}ngen von maximal
200\,Basen beschr\"{a}nkt, da die Kupplungseffizienz mit jeder Verl\"{a}ngerungsrunde
abnimmt und ab 200 Zyklen typisch unter 98\,\% sinkt, was zu einem
exponentiell wachsenden Anteil fehlerhafter Sequenzen f\"{u}hrt.
Template-unabh\"{a}ngige Enzyme --- insbesondere die
\emph{Terminale Desoxy-nukleotidyltransferase~(TdT)} --- \"{u}berwinden
diese Grenze grunds\"{a}tzlich, da sie Nukleotide enzymatisch anf\"{u}gen
ohne strukturellen Verfall des Substrats.
Aktuelle TdT-basierte Systeme (Nuclera Nucleomics, Ansa Biotechnologies)
erreichen zuverl\"{a}ssig 500-mere; langfristig sind Synthesel\"{a}ngen
bis 1000+ Basen realisierbar.

F\"{u}r das DSS-System hat dies direkte, formal quantifizierbare Auswirkungen.
Nach Satz~\ref{satz:max_dichte} w\"{a}chst die maximale Informationsdichte
linear mit~$k$:
\[
    I_{\max}(k, r_{\mathrm{rs}}, \rho) \;=\;
    k \cdot (2 - r_{\mathrm{rs}}) \cdot \frac{1}{1-\rho}.
\]
F\"{u}r $k = 1000$ (statt heute 150\,Basen) bei gleichem
$r_{\mathrm{rs}} = 0{,}13$ und $\rho = 0{,}8$ ergibt sich:
\[
    I_{\max}(1000,\, 0{,}13,\, 0{,}8) = 1000 \cdot 1{,}87 \cdot 5
    = 9350 \;\text{Bit/Oligo} = 1168 \;\text{Byte/Oligo}.
\]
Dies entspricht einer Steigerung um den Faktor $6{,}67$ gegen\"{u}ber dem
150-mer-System.
Zus\"{a}tzlich w\"{a}chst mit der Sequenzl\"{a}nge $k$ die Anzahl der
Watson-Crick-Kanten im Datengraphen $G_D$ quadratisch anwachsend an
(jede Base kann mit bis zu $k-1$ Gegenbases eine komplementäre
Kante eingehen), was die erreichbare Komprimierungsrate $\rho$ tendenziell
erh\"{o}ht und damit den Effekt \"{u}berproportional verst\"{a}rkt.
Eine formale Analyse dieser R\"{u}ckkopplung zwischen Oligomerl\"{a}nge
und Komprimierungsrate ist Gegenstand laufender Forschung.

%---
\subsubsection*{3.\; In-Memory-DNA: Hydrogel-Arrays mit parallellem Nanoporen-Zugriff}

DNA-vernetzte Hydrogel-Matrices erlauben die r\"{a}umliche Fixierung einzelner
DNA-Sequenz\-bereiche ohne Extraktion aus dem Speichermedium.
Der Signatur-Index $\mathcal{I}$ (Definition~\ref{def:index}) kann dabei
als r\"{a}umliche Adressierstruktur \"{u}ber das Gel-Gitter dienen:
jeder Gel-Knoten mit einem bestimmten Raumindex $(x,y,z)$ enth\"{a}lt
gezielt die Oligomere mit Signatur $\boldsymbol{\sigma}^D = f(x,y,z)$.

Sei $N_{\mathrm{pore}}$ die Anzahl paralleler Nanoporen-Kan\"{a}le.
Die Zugriffszeit f\"{u}r einen beliebigen Block $D_i$ im Pool der Gr\"{o}\ss{}e~$M$
betr\"{a}gt klassisch (ohne Parallelisierung):
\[
    t_{\mathrm{access}}^{\mathrm{seq}} \;=\;
    t_{\mathrm{PCR}} + t_{\mathrm{seq}} \cdot \frac{M}{c_{\mathrm{cover}}},
\]
wobei $c_{\mathrm{cover}} \ge 50$ die erforderliche \"{U}berdeckungstiefe ist.
F\"{u}r $M = 10^6$ und $t_{\mathrm{seq}} = 10$\,ms/Oligo:
$t_{\mathrm{access}}^{\mathrm{seq}} \approx 3$ Stunden.

Mit $N_{\mathrm{pore}}$ parallelen Nanoporen und Index-gesteuerter Selektion:
\[
    t_{\mathrm{access}}^{\mathrm{parallel}} \;=\;
    O\!\left(\frac{\log M \cdot t_{\mathrm{seq}}}{N_{\mathrm{pore}}}\right).
\]
Oxford Nanopore PromethION bietet bereits 2675 parallele Kan\"{a}le;
n\"{a}chste Generationen zielen auf $N_{\mathrm{pore}} > 10^5$ ab,
was Zugriffszeiten unter einer Sekunde f\"{u}r Pools bis $10^9$ Oligomere
erm\"{o}glicht.
Die Kopplung mit dem Subgraph-Signatur-Index erm\"{o}glicht dar\"{u}ber hinaus
\emph{semantischen Direktzugriff}: Statt einer Adresse wird eine strukturelle
Anfrage $G_Q$ gestellt, die alle Bl\"{o}cke mit
$G_{D_i} \preceq_{\mathrm{rot}} G_Q$ findet ---
ein direkter Subgraph-Join im physikalischen Speichermedium.

%---
\subsubsection*{4.\; CRISPR-basiertes gezieltes Schreiben: echter Random-Write in DNA}

Bisher gilt DNA-Speicher als \emph{Write-Once}-Medium: einmal synthetisiert,
nicht modifizierbar ohne Re-Synthese. CRISPR-Cas-basierte Editierwerkzeuge
brechen dieses Paradigma fundamental:
\begin{itemize}
    \item \textbf{Baseeditoren (BE3, ABEmax, CBE):} Erlauben Konversionen
          $A \to G$ und $C \to T$ (bzw.\ $C \to T$ und $G \to A$) ohne
          DNA-Doppelstrangbruch, mit Pr\"{a}zision von~$> 99\,\%$ f\"{u}r
          die Zielbase.
    \item \textbf{Prime Editing (PE3, PE3b):} Cas9-Nickase verbunden mit
          Reverser Transkriptase erm\"{o}glicht beliebige Substitutionen,
          Insertionen und Deletionen in einem einzigen Editing-Schritt.
          Zielbereichsgr\"{o}\ss{}e: bis 80 Basen.
    \item \textbf{CRISPR-Cas12a (Cpf1):} Erzeugt klebrige Enden f\"{u}r
          pr\"{a}zise Insertion gr\"{o}\ss{}erer Sequenzst\"{u}cke
          ($\le 500$\,bp).
\end{itemize}

Im DSS-Formalismus entspricht ein einzelner CRISPR-Edit einer lokalen
Modifikation des DNA-Datengraphen $G_D$.
Sei $G_D' = G_D \oplus \Delta$ das editierte Ergebnis, wobei
$\Delta$ die Kantenver\"{a}nderung (Addition oder Entfernung einer
Watson-Crick-Kante infolge der Basissubstitution) beschreibt.
Die neue Signatur $\boldsymbol{\sigma}^{D'}$ weicht von $\boldsymbol{\sigma}^D$
in genau $j$ Komponenten ab, wobei $j \le 2k$ gilt
(jede editierte Base ist in h\"{o}chstens $k$ Watson-Crick-Kanten involviert),
bei einer Einzelbaseneditierung jedoch praktisch $j = O(1)$.

Der Subgraph-Index $\mathcal{I}$ muss nach einem Edit lediglich in $O(j)$
Eintr\"{a}gen aktualisiert werden --- eine Verbesserung gegen\"{u}ber
naiver Re-Synthese des gesamten Oligos um den Faktor~$k/j$.
Damit w\"{a}re das DSS das erste DNA-Speichersystem mit echtem
$O(1)$-Random-Write pro Bit-\"{A}nderung im physikalischen Medium.

Offene Herausforderungen: \emph{(i)}~Zuverl\"{a}ssige Zuleitung der CRISPR/Cas-Komplexe
zu spezifischen Oligomeren in einem dichten Pool ohne Off-Target-Aktivit\"{a}t;
\emph{(ii)}~Skalierbarkeit auf $M > 10^{12}$ Oligomere;
\emph{(iii)}~Formale Fehleranalyse des Edits und seine Integration in den
Reed-Solomon-Verifikationsschritt (Satz~\ref{satz:retrieval}, Schritt~7).
Vielversprechende Fortschritte in der Nanoporen-gesteuerten CRISPR-Zustellung
und in Single-Molecule-Monitoring machen einen ersten Laborbeweis im
Zeithorizont von 5--10 Jahren realistisch.

%---
\subsubsection*{5.\; Maschinelles Lernen f\"{u}r Kodierung, Basecalling und Komprimierung}

Neuronale Netze haben das Basecalling revolutioniert~\cite{deamer2016}.
F\"{u}r das DSS ergeben sich folgende ML-Anwendungsfelder:

\begin{enumerate}
    \item \textbf{Gelernte Signatur-Hashfunktionen.}
          Anstelle der algebraischen Signaturfunktion
          (Definition~\ref{def:signatur}) k\"{o}nnte ein trainiertes
          Einbettungsmodell $f_\theta: \Sigma^k \to \mathbb{R}^d$ eine
          \emph{semantische Einbettung} lernen, die strukturell
          \"{a}hnliche Sequenzen nah im $\mathbb{R}^d$ platziert.
          Dies entspricht einem Approximate Nearest Neighbor~(ANN)-Problem.
          Locality-Sensitive Hashing~(LSH) \"{u}ber Subgraph-Signaturen
          bietet $O(1)$ Lookup bei $O(k)$ Hash-Berechnung, womit der
          exakte $O(\log M)$-Index f\"{u}r gro\ss{}e Pools weiter beschleunigt
          werden kann.
    \item \textbf{Graph Neural Networks f\"{u}r approximatives Subgraph-Matching.}
          Unter SETH ist $O(k^3)$ f\"{u}r den exakten Subgraph-Test asymptotisch
          optimal (Satz~\ref{satz:optimalitaet_gesamt}).
          F\"{u}r approximatives Matching hingegen kann ein Graph Neural
          Network~(GNN) den Test in $O(k \cdot d_{\mathrm{GNN}})$ durchf\"{u}hren,
          wobei $d_{\mathrm{GNN}}$ die Tiefe des Netzes ist.
          Bei $k = 150$, $d_{\mathrm{GNN}} = 4$ ergibt sich eine Beschleunigung
          um Faktor $k^2/d_{\mathrm{GNN}} \approx 5600$.
          Im Komprimierungsschritt wird approximatives Subgraph-Matching
          (mit kontrollierter Fehlerrate $\beta > 0$) toleriert: ein Block $D_i$
          wird als redundant markiert, wenn das GNN-Modell
          $\Pr[G_{D_i} \preceq G_{D_j}] > 1 - \beta$ f\"{u}r ein
          $D_j \in \mathcal{M}$ ausgibt.
    \item \textbf{Neuronale GC-Balance-Optimierung.}
          Die aktuelle GC-Balance-Pr\"{u}fung (Schreib-Algorithmus, Schritt~3)
          ist regelbasiert und heuristisch. Ein Reinforcement-Learning-Agent
          k\"{o}nnte Kodierungsparameter adaptiv anpassen, um GC-Homogenit\"{a}t
          \"{u}ber alle Oligos des Pools zu maximieren, was die Hybridisierungsrate
          beim Sequenzieren stabilisiert und damit die effektive Fehlerrate
          $\varepsilon$ in Satz~\ref{satz:kapazitaet} senkt.
\end{enumerate}

%---
\subsubsection*{6.\; Zweistufige Fehlerkorrektur: RS + Signatur-ECC}

Anstelle einer vollst\"{a}ndigen Quantenfehlerkorrektur, die DNA-Qubits
erfordern w\"{u}rde, l\"{a}sst sich ein konkret realisierbares
\emph{zweistufiges ECC-Schema} konstruieren, das die klassische
Reed-Solomon-Schicht um eine Signatur-ECC-Schicht erg\"{a}nzt:

\begin{enumerate}
    \item \textbf{Stufe~1 (Sequenz-ECC):} Reed-Solomon $\mathrm{RS}(n_{\mathrm{rs}},k_{\mathrm{rs}})$
          auf Basisebene (Definition~\ref{def:reedsolomon}), korrigiert
          bis zu $t$ Substitutionsfehler pro Oligo.
    \item \textbf{Stufe~2 (Signatur-ECC):} BCH- oder LDPC-Code
          \"{u}ber dem Signaturvektor $\boldsymbol{\sigma}^D \in \mathbb{N}^k$.
          Da $\sigma_j \in \mathbb{N}$ eine algebraische Struktur besitzt
          (ganzzahlige Codierung der Spaltensummen), ist ein linearer
          Block-Code in $\mathbb{F}_2^{\lceil\log_2\sigma_{\max}\rceil}$
          direkt anwendbar.
          Dieser detektiert Burst-Fehler, die Stufe~1 \"{u}berschreiten,
          und markiert sie als \emph{Erasure} statt als stille Datenkorruption,
          was der Subgraph-Verifikationsschritt (Algorithmus~\ref{satz:retrieval},
          Schritt~7) auswertet.
\end{enumerate}

Die residuale unkorrekte-Bit-Rate ergibt sich als Produkt beider
Fehlerwahrscheinlichkeiten:
\[
    P_{\mathrm{err}}^{(2)} \;\approx\;
    P_{\mathrm{err}}^{\mathrm{RS}} \cdot P_{\mathrm{miss}}^{\mathrm{Sig-ECC}}
    \;\ll\; P_{\mathrm{err}}^{\mathrm{RS}}.
\]
F\"{u}r $k=150$, $t=2$, $\varepsilon = 10^{-3}$ ergibt sich
$P_{\mathrm{err}}^{\mathrm{RS}} \approx 10^{-14}$;
mit einem BCH$(31,16)$-Code f\"{u}r die Signatur-Schicht reduziert sich
dies auf $P_{\mathrm{err}}^{(2)} < 10^{-26}$ ---
eine Gr\"{o}\ss{}enordnung, die selbst f\"{u}r sicherheitskritische
Millionen-Jahres-Archivierung ausreicht
(Satz~\ref{satz:haltbarkeit}, $t^* > 10^6$ Jahre).

%---
\subsubsection*{7.\; Zellbasierte DNA-Speicherung (\textit{in vivo})}

Lebende Zellen k\"{o}nnen als biochemische Speichermedien dienen:
In \emph{E.\,coli} wurden bereits Bin\"{a}rdaten von $\sim\!100$\,kB
in das chromosomale Bakteriengenom integriert
(mit CRISPR-gesteuerten Rekombinationsmodulen).
Das DSS-Modell l\"{a}sst sich grunds\"{a}tzlich auf \emph{lebende} Organismen
ausweiten:

Im in-vivo-Kontext entspr\"{a}che der DNA-Datengraph $G_D$ einer
Erweiterung um epige\-netische Kanten---Methylierungsmuster
(CpG-Methylierung, $m^5C$, $m^6A$) kodieren zus\"{a}tzlich 1\,Bit/Base.
Der erweiterte Datengraph:
\[
    G_D^{\mathrm{epi}} \;=\;
    (V_D,\; E_{\mathrm{back}} \cup E_{\mathrm{wc}} \cup E_{\mathrm{meth}})
\]
k\"{o}nnte die Informationsdichte von $2$\,Bit/Base theoretisch auf
$3$\,Bit/Base steigern --- ein Gewinn von 50\,\%.
Die Subgraph-Signaturfunktion~(Definition~\ref{def:signatur}) muss dabei
auf den erweiterten Graphen $G_D^{\mathrm{epi}}$ verallgemeinert werden,
indem Kantengewichte $w_e \in \{0,1\}$ (methyliert / nicht methyliert)
in die Spalten-Signatur einflie\ss{}en:
\[
    \sigma_j^{\mathrm{epi}} \;=\; \sum_{i=0}^{k-1}
    (A^D_{ij} + w_{ij} \cdot 2^k) \cdot 2^i \;+\; j \cdot 2^{2k}.
\]
Herausforderungen: Methylierungs\-muster sind zellteilungsabh\"{a}ngig
und damit zeitlich instabil. Formale Fehlerkorrekturschemata f\"{u}r diesen
epigenetischen Kanal sowie die Integration in die Reed-Solomon-Schicht bleiben
weitgehend offene Forschungsfragen.

%---
\subsubsection*{8.\; Standardisierung und Interoperabilit\"{a}t}

Die technische Reife des DNA-Speichers macht die Entwicklung von
Interoperabilit\"{a}tsstan\-dards notwendig.
Das ISO Technical Committee~276 (Biotechnology, Synthetic Biology)
und das IETF-Consortium f\"{u}r biologische Datenformate haben erste
Entw\"{u}rfe vorgelegt.
Das DSS-System kann als \emph{Referenzarchitektur} f\"{u}r DNA-Speicher\-protokolle
dienen, da es das erste vollst\"{a}ndig formal spezifizierte und bewiesene
System seiner Art ist.

Ein standardisiertes DSS-Protokoll w\"{u}rde folgende Komponenten umfassen:
\begin{itemize}
    \item \textbf{DNA-Dateiformat (DFF):}
          Prim\"{a}re Sequenz + Pr\"{a}ambel
          (Block-ID, Versionsnummer, ECC-Typ, GC-Inversions-Flag)
          als standardisiertes $5'$-$3'$-Layout.
    \item \textbf{Signaturprotokoll:}
          Kanonische Berechnung von $\boldsymbol{\sigma}^D$
          (Definition~\ref{def:signatur}) als plattformunabh\"{a}ngige
          Referenz\-implementierung, MIME-Typ \texttt{application/x-dna-dss}.
    \item \textbf{Fehlerkorrektur-Profil:}
          Standardisierte RS-Parameter $(n_{\mathrm{rs}}, k_{\mathrm{rs}}, t)$
          f\"{u}r drei Klassen:
          \emph{High-Density} ($t=1$),
          \emph{High-Durability} ($t=4$),
          \emph{Ultra-Reliable} ($t=8$, zweistufige ECC).
    \item \textbf{Inter-Pool-Protokoll:}
          HTTP/2-\"{a}hnliches Request-Response-Schema mit Signatur-basierter
          Adressierung f\"{u}r den Zugriff auf verteilte DNA-Pools in
          geographisch-verteilten Rechenzentren.
\end{itemize}

Die formalen Korrektheitseigenschaften des DSS
(S\"{a}tze~\ref{satz:kodierung}--\ref{satz:optimalitaet_gesamt})
bilden dabei die unveranderte mathematische Spezifikation des Standards ---
analog zur Rolle der RFC-Dokumente f\"{u}r das TCP/IP-Protokoll.

%---
\subsubsection*{9.\; Wirtschaftliche Perspektive und Technologie-Roadmap}

Die \"{o}konomische Wettbewerbsf\"{a}higkeit von DNA-Speichern h\"{a}ngt
prim\"{a}r von den Synthese- und Sequenzierkosten ab.
Tabelle~\ref{tab:kosten_roadmap} zeigt eine konservative Projektion.

\begin{table}[H]
\centering
\caption{Kostenprojektionen und DSS-Break-Even-Analyse 2024--2035.
$C_{\mathrm{syn}}$: Synthesekosten [\$/Base],
$C_{\mathrm{seq}}$: Sequenzierkosten [\$/Gb],
$C_{\mathrm{eff}}$: effektive Synthesekosten nach DSS-Komprimierung
($\rho = 0{,}8$, Faktor~5 Reduktion).
Vergleichswert: Tape-Speicher $\approx \$10^3$/EB (2024).}
\label{tab:kosten_roadmap}
\begin{tabular}{lrrrrl}
\toprule
Jahr & $C_{\mathrm{syn}}$ & $C_{\mathrm{seq}}$ &
$C_{\mathrm{eff}}$  & \$/EB & Wettbewerb \\
\midrule
2024 & $10^{-1}$           & $10^{-3}$          & $2{\times}10^{-2}$  & $10^{18}$ & HDD $10^4$ \\
2026 & $10^{-2}$           & $10^{-4}$          & $2{\times}10^{-3}$  & $10^{16}$ & Tape $10^3$ \\
2028 & $10^{-3}$           & $10^{-5}$          & $2{\times}10^{-4}$  & $10^{14}$ & Tape $10^3$ \\
2030 & $5{\times}10^{-4}$  & $5{\times}10^{-6}$ & $10^{-4}$           & $5{\times}10^{12}$ & Tape $10^3$ \\
2035 & $10^{-5}$           & $10^{-7}$          & $2{\times}10^{-6}$  & $10^9$    & Tape $10^3$ \\
\bottomrule
\end{tabular}
\end{table}

Der Break-Even mit Tape-Speicher (aktuell $\sim\!\$1000$/EB) ist nach
dieser Projektion zwischen 2030 und 2035 zu erwarten.
Der Subgraph-Komprimierungsfaktor $(1-\rho)^{-1} = 5$ (bei $\rho=0{,}8$)
reduziert die Synthesekosten direkt um $80\,\%$, was den Break-Even
um sch\"{a}tzungsweise 3--5 Jahre vorverlagert.
Dies unterstreicht den \"{o}konomischen Wert der Subgraph-Komprimierung
unabh\"{a}ngig von ihrer theoretischen Eleganz.

%---
\subsubsection*{10.\; Globales DNA-Archiv der Menschheit (GenVault)}

Die formalen Garantien dieser Arbeit legen die mathematische Grundlage
f\"{u}r ein auf Jahrtausende ausgelegtes, globales Digitaldaten-Archiv.
Nach Satz~\ref{satz:haltbarkeit} betr\"{a}gt die garantierte Lebensdauer
des DSS bei 80\,\% Redundanz und kryogenischer Lagerung
($-196$\,\textdegree{}C, fl\"{u}ssiger Stickstoff,
$\lambda \approx 10^{-9}$/Jahr nach Arrhenius-Extrapolation):
\[
    t^* \;=\; \frac{-\ln(1-0{,}8)}{10^{-9}/\mathrm{Jahr}}
    \;\approx\; 1{,}6 \times 10^9 \;\text{Jahre}.
\]
Dies \"{u}bersteigt das gesch\"{a}tzte zuk\"{u}nftige Alter der Erde um
eine Milliarde Jahre.

Ein realistisches Archiv-Szenario:
\begin{itemize}
    \item \textbf{Kapazit\"{a}t:} 200\,ZB (gesamte projizierte Daten\-produktion bis 2040,
          nach IDC Global DataSphere).
    \item \textbf{DNA-Dichte:} $2{,}15 \times 10^{15}$\,B/g (nativ)
          $\times\,5$ (DSS-Komprimierung, $\rho=0{,}8$)
          $= 1{,}075 \times 10^{16}$\,B/g.
    \item \textbf{Ben\"{o}tigte Masse:}
          $\frac{200 \times 10^{21}\,\text{B}}{1{,}075 \times 10^{16}\,\text{B/g}}
          \approx 18{,}6 \times 10^6\,\text{g} = 18{,}6\,\text{Tonnen}$
          lyophilisierte DNA.
    \item \textbf{Volumen} (Dichte $\approx 1{,}7\,\text{g/cm}^3$):
          $\approx 11\,\text{m}^3$ --- weniger als ein Seefrachtcontainer.
    \item \textbf{Energiebedarf:} Kryolagerung bei $-20$\,\textdegree{}C,
          dauerhaft $\approx 1$\,kW elektrischer Leistung.
    \item \textbf{Vergleich:} Ein modernes Hyperscale-Rechenzentrum
          (z.\,B.\ Microsoft Azure Region) belegt $> 10^5\,\text{m}^3$
          und verbraucht $\approx 100$\,MW Dauerleistung f\"{u}r
          $\approx 10$\,ZB Speicher.
\end{itemize}

Das DSS-Archiv w\"{u}rde also den 20-fachen Speicherinhalt eines Hyperscale-Zentrums
in einem zehntausendmal kleineren Volumen bei einem Zehnmillionstel des
Energieaufwands bereitstellen --- und dies f\"{u}r $> 10^9$ Jahre.

%---
\subsubsection*{Offene mathematische Probleme}

Abschlie\ss{}end werden die wichtigsten offenen Fragen f\"{u}r zuk\"{u}nftige
Forschung benannt:
\begin{enumerate}
    \item \textbf{Untere Schranke f\"{u}r erreichbare Komprimierung.}
          Wie gro\ss{} ist der minimale Redundanzgrad $\rho_{\min}$,
          der f\"{u}r eine gegebene Verteilung $\mathcal{D}$ \"{u}ber
          Datenbl\"{o}cke erreichbar ist?
          Dies h\"{a}ngt mit der Kolmogorov-Komplexit\"{a}t der Eingabe
          zusammen und ist im Allgemeinen nicht berechenbar.
          F\"{u}r parameterierte Familien~$\mathcal{D}$ sind scharfe
          Schranken jedoch m\"{o}glich und bilden eine offene Frage
          der algorithmischen Informationstheorie.
    \item \textbf{Komplexit\"{a}t des maximalen Kerngenoms.}
          Die Berechnung des optimalen maximalen Kerngenoms $\mathcal{M}$
          unter der $\preceq_{\mathrm{rot}}$-Ordnung ist m\"{o}glicherweise
          NP-schwer (Reduktion auf Maximum-Weight-Clique in einem
          Kompatibilit\"{a}tsgraphen).
          Formale Komplexit\"{a}tsklassifikation, Inapproximierbarkeitsresultate
          und PTAS-Schemata stehen aus.
    \item \textbf{Online-Subgraph-Komprimierung.}
          Der aktuelle Algorithmus ist batch-orientiert.
          Ein Online-Algorithmus, der Bl\"{o}cke sequenziell verarbeitet
          und dabei eine $\alpha$-Approxi\-mation des optimalen Kerngenoms
          aufrechterh\"{a}lt, fehlt bisher.
          Verbindungen zu Streaming-Algorithmen f\"{u}r Graph-Probleme
          sind zu erwarten.
    \item \textbf{Topologische DNA-Adressierung.}
          K\"{o}nnen dreidimensionale DNA-Origami-Struk\-turen Datengraphen
          $G_D$ entsprechen, deren Subgraph-Signatur geometrische
          Invarianten kodiert (Verschlingungszahl $lk$, Genus $g$)?
          Dies k\"{o}nnte zu einer \emph{topologischen DNA-Adressierung}
          f\"{u}hren, bei der strukturelle topologische Eigenschaften
          als Speicheradresse dienen --- ein v\"{o}llig neuartiges
          Adressierungsparadigma.
    \item \textbf{Quantitative Theorie der GC-Ausgewogenheit.}
          F\"{u}r welche Wahrscheinlichkeitsverteilungen $\mathcal{P}$
          \"{u}ber $\mathcal{B}^{2k}$ gilt
          $\mathbb{E}_{D \sim \mathcal{P}}[|GC(\phi(D))|/k] = 1/2$?
          Eine vollst\"{a}ndige Charakterisierung w\"{u}rde den
          Inversionsschritt (Schreib-Algorithmus, Schritt~3) durch eine
          optimale Bijektion $\phi^*$ ersetzen.
\end{enumerate}

\subsection{Implementierungshinweis}

Die Implementierung des Subgraph Algorithmus ist verfügbar unter:
\url{https://gitlab.com/epp-group/subgraph}.

\begin{thebibliography}{12}

    \bibitem{epp2026}
        Stephan Epp.
        \textit{Der Subgraph Algorithmus}. 2026.
        \url{https://gitlab.com/epp-group/subgraph}.

    \bibitem{epp_dna2026}
        Stephan Epp.
        \textit{Graphenbasierte DNA-Sequenzierung mittels des Subgraph
        Algorithmus}.
        Universit\"{a}t Bielefeld, April 2026.

    \bibitem{church2012}
        G.\,M. Church, Y. Gao und S. Kosuri.
        Next-Generation Digital Information Storage in DNA.
        \textit{Science}, 337(6102):1628, 2012.

    \bibitem{goldman2013}
        N. Goldman et al.
        Towards practical, high-capacity, low-maintenance information
        storage in synthesized DNA.
        \textit{Nature}, 494:77--80, 2013.

    \bibitem{organick2019}
        L. Organick et al.
        Random access in large-scale DNA data storage.
        \textit{Nature Biotechnology}, 36:242--248, 2018.

    \bibitem{reedsolomon1960}
        I.\,S. Reed und G. Solomon.
        Polynomial Codes Over Certain Finite Fields.
        \textit{J. SIAM}, 8(2):300--304, 1960.

    \bibitem{shannon1948}
        C.\,E. Shannon.
        A Mathematical Theory of Communication.
        \textit{Bell System Technical Journal}, 27:379--423, 1948.

    \bibitem{watson1953}
        J.\,D. Watson und F.\,H.\,C. Crick.
        Molecular Structure of Nucleic Acids.
        \textit{Nature}, 171:737--738, 1953.

    \bibitem{abboud2015}
        A. Abboud, A. Backurs und V. Vassilevska Williams.
        Tight Hardness Results for LCS and Other Sequence Similarity
        Measures.
        \textit{FOCS 2015}, 2015.

    \bibitem{erlich2017}
        Y. Erlich und D. Zielinski.
        DNA Fountain enables a robust and efficient storage architecture.
        \textit{Science}, 355(6328):950--954, 2017.

    \bibitem{deamer2016}
        D. Deamer, M. Akeson und D. Branton.
        Three decades of nanopore sequencing.
        \textit{Nature Biotechnology}, 34:518--524, 2016.

    \bibitem{ceze2019}
        L. Ceze, J. Nivala und K. Strauss.
        Molecular digital data storage using DNA.
        \textit{Nature Reviews Genetics}, 20:456--466, 2019.

\end{thebibliography}

\end{document}
